% Annexe A -- Rappels d'analyse. A inclure apres \appendix par % Annexe A -- Rappels d'analyse. A inclure apres \appendix par % Annexe A -- Rappels d'analyse. A inclure apres \appendix par % Annexe A -- Rappels d'analyse. A inclure apres \appendix par \input{appA_analyse.tex}.
% Etiquettes prefixees appA: ; references bibliographiques prefixees app: (refs_app.tex).
\chapter{Rappels d'analyse}\label{appA:chap}

Cette annexe rassemble les outils d'analyse que les chapitres utilisent sans toujours les redémontrer :
\begin{itemize}
\item l'interversion d'une limite et d'une intégrale ou d'une série (§\ref{appA:sec:param}) ;
\item les transformées de Fourier et de Laplace (§\ref{appA:sec:fourier}) ;
\item la formule sommatoire de Poisson, dans le cas général (§\ref{appA:sec:poisson}) ;
\item les fonctions spéciales rencontrées : $\Gamma$, $\operatorname{erfc}$, $J_0$, $K_0$, $Y_0$ et la fonction de
Struve $\mathbf H_0$ (§\ref{appA:sec:speciales}) ;
\item les réseaux du plan et le domaine fondamental de $SL_2(\Z)$ (§\ref{appA:sec:sl2}).
\end{itemize}
Le niveau visé est celui de la deuxième année de classe préparatoire. L'intégrale est celle de Lebesgue, mais il suffit
de savoir appliquer les deux théorèmes de convergence du programme.

\begin{warnbox}[: niveaux des énoncés de cette annexe]
Les démonstrations écrites ici sont des démonstrations humaines \tierC{} : aucune n'est vérifiée par Lean dans ce livre.
Les énoncés admis (convergence dominée, théorème de Fejér, propriétés des fonctions de Bessel et de Struve tirées du
DLMF~\cite{app:dlmf}) sont de niveau \tierL. La Mathlib épinglée contient des versions formelles de plusieurs d'entre
eux : la convergence dominée (\lean{MeasureTheory.tendsto_integral_of_dominated_convergence}, fichier
\texttt{MeasureTheory/Integral/DominatedConvergence.lean}) et la formule de Poisson (fichiers
\texttt{Analysis/Fourier/PoissonSummation.lean} et \texttt{Analysis/SpecialFunctions/Gaussian/PoissonSummation.lean}).
Nous avons seulement constaté que ces fichiers existent ; nous n'avons ni comparé leurs énoncés à ceux de cette annexe,
ni relié ces énoncés aux démonstrations ci-dessous.
\end{warnbox}

%=====================================================================================================================
\section{Séries et intégrales dépendant d'un paramètre}\label{appA:sec:param}

\subsection{Le théorème de convergence dominée}

\begin{theorem}[Convergence dominée \tierL]\label{appA:thm:cvd}
Soit $I$ un intervalle de $\R$ et $(f_n)$ une suite de fonctions continues par morceaux sur $I$. On suppose que :
\begin{enumerate}
\item $(f_n)$ converge simplement sur $I$ vers une fonction $f$ continue par morceaux ;
\item il existe $\varphi:I\to\R_+$ intégrable sur $I$ telle que $|f_n(x)|\le\varphi(x)$ pour tout $n$ et tout $x$
(hypothèse de domination).
\end{enumerate}
Alors les $f_n$ et $f$ sont intégrables sur $I$, et $\int_If_n\to\int_If$.
\end{theorem}

Le point essentiel est que la fonction dominante $\varphi$ ne dépend pas de $n$. Sans elle, la conclusion peut échouer :
$f_n=n\,\mathbf 1_{]0,1/n]}$ tend simplement vers $0$ sur $]0,1]$, alors que $\int f_n=1$ pour tout $n$.

\begin{theorem}[Interversion série--intégrale \tierL]\label{appA:thm:serieint}
Soit $(u_n)$ une suite de fonctions intégrables sur $I$ telle que $\sum_n u_n$ converge simplement vers une fonction
continue par morceaux et $\sum_n\int_I|u_n|<\infty$. Alors $\sum_nu_n$ est intégrable et
$\int_I\sum_nu_n=\sum_n\int_Iu_n$.
\end{theorem}

\subsection{Continuité et dérivation sous l'intégrale}

\begin{theorem}[Continuité sous l'intégrale]\label{appA:thm:cont}
Soit $A$ une partie de $\R$ et $f:A\times I\to\C$ telle que $f(\cdot,x)$ est continue pour tout $x$, $f(a,\cdot)$ est
continue par morceaux pour tout $a$, et $|f(a,x)|\le\varphi(x)$ avec $\varphi$ intégrable. Alors
$F(a)=\int_If(a,x)\,dx$ est continue sur $A$.
\end{theorem}

\begin{humanproof}
Soit $a_n\to a$ dans $A$. Les fonctions $x\mapsto f(a_n,x)$ convergent simplement vers $f(a,\cdot)$ et sont dominées par
$\varphi$. Le théorème~\ref{appA:thm:cvd} donne $F(a_n)\to F(a)$, et la caractérisation séquentielle conclut.
\end{humanproof}

\begin{theorem}[Dérivation sous l'intégrale]\label{appA:thm:deriv}
Soit $J$ un intervalle et $f:J\times I\to\C$. On suppose que $f(a,\cdot)$ est intégrable pour tout $a\in J$, que
$\partial f/\partial a$ existe, est continue en $a$ et continue par morceaux en $x$, et qu'il existe $\psi$ intégrable
avec $|\partial_af(a,x)|\le\psi(x)$ pour tous $a,x$. Alors $F(a)=\int_If(a,x)\,dx$ est de classe $\mathcal C^1$ sur $J$,
et $F'(a)=\int_I\partial_af(a,x)\,dx$.
\end{theorem}

\begin{humanproof}
Fixons $a\in J$ et une suite $h_n\to0$, $h_n\ne0$, avec $a+h_n\in J$. Posons
$g_n(x)=\bigl(f(a+h_n,x)-f(a,x)\bigr)/h_n$. Par définition de la dérivée partielle, $g_n(x)\to\partial_af(a,x)$ pour
tout $x$. Par l'inégalité des accroissements finis appliquée à $t\mapsto f(t,x)$ entre $a$ et $a+h_n$,
$|g_n(x)|\le\sup_t|\partial_af(t,x)|\le\psi(x)$. Le théorème~\ref{appA:thm:cvd} donne
$\bigl(F(a+h_n)-F(a)\bigr)/h_n=\int g_n\to\int\partial_af(a,x)\,dx$. La continuité de $F'$ découle du
théorème~\ref{appA:thm:cont} appliqué à $\partial_af$.
\end{humanproof}

\begin{remark}[Domination locale]
Il suffit en pratique que la domination ait lieu sur tout segment $[\alpha,\beta]\subset J$ : la dérivabilité est une
propriété locale. C'est ainsi qu'on traite les transformées de Laplace ci-dessous, où la domination n'est uniforme que
pour $t\ge\alpha>0$.
\end{remark}

\begin{example}[Moments gaussiens]\label{appA:ex:moments}
Pour $\alpha>0$, $G(\alpha)=\int_\R e^{-\pi\alpha x^2}dx=\alpha^{-1/2}$ (intégrale de Gauss). Sur $[\alpha_0,\infty[$ avec
$\alpha_0>0$, $|\partial_\alpha e^{-\pi\alpha x^2}|=\pi x^2e^{-\pi\alpha x^2}\le\pi x^2e^{-\pi\alpha_0x^2}$, qui est
intégrable. Le théorème~\ref{appA:thm:deriv} donne $\int\pi x^2e^{-\pi\alpha x^2}dx=-G'(\alpha)=\tfrac12\alpha^{-3/2}$.
En dérivant $m$ fois, $\int(\pi x^2)^me^{-\pi\alpha x^2}dx=\frac{(2m)!}{4^mm!}\,\alpha^{-1/2-m}$ (exercice~\ref{appA:ex:moments}).
C'est la même technique qui, au chapitre~1, calcule la transformée de Fourier d'une gaussienne
(théorème~\ref{ch1:thm:poisson}, étape~1).
\end{example}

\subsection{Séries de fonctions}

\begin{proposition}[Convergence normale \tierL]\label{appA:prop:normale}
Soit $(u_n)$ des fonctions continues sur une partie $A$ de $\R^d$ avec $\sum_n\sup_A|u_n|<\infty$. Alors
$\sum_nu_n$ converge uniformément sur $A$ et sa somme est continue. Si de plus les $u_n$ sont $\mathcal C^1$ sur un
ouvert et $\sum_n\sup|\nabla u_n|<\infty$, la somme est $\mathcal C^1$ et se dérive terme à terme.
\end{proposition}

\begin{example}[Séries thêta]
Pour un réseau $\Lambda$ du plan et $\alpha\ge\alpha_0>0$, on a $|x|^2\ge c\,(m^2+n^2)$ pour $x=mu+nv$ (chapitre~1),
donc $e^{-\pi\alpha|x|^2}\le e^{-\pi\alpha_0c(m^2+n^2)}$, terme général d'une série convergente. La série
$\Theta_\Lambda(\alpha)$ converge donc normalement sur $[\alpha_0,\infty[$, et elle y est $\mathcal C^\infty$ en
$\alpha$ : on dérive terme à terme autant qu'on veut. Les sommes de réseau $S(\alpha)=\sum|P|^2e^{-\pi\alpha|P|^2}$ du
chapitre~3 sont précisément $-\Theta_\Lambda'(\alpha)/\pi$.
\end{example}

%=====================================================================================================================
\section{Transformées de Fourier et de Laplace}\label{appA:sec:fourier}

\subsection{La transformée de Fourier}

\begin{definition}
Pour $f:\R^d\to\C$ intégrable, la \emph{transformée de Fourier} de $f$ est
\[\hat f(\xi)=\int_{\R^d}f(x)\,e^{-2\pi i\,\xi\cdot x}\,dx,\qquad\xi\in\R^d .\]
\end{definition}
C'est la convention du chapitre~1 : le facteur $2\pi$ est dans l'exponentielle, et la formule d'inversion n'a alors
aucun facteur devant l'intégrale. D'autres livres, en physique surtout, écrivent $\int f(x)e^{-ik\cdot x}dx$ ; on passe
de l'une à l'autre par $k=2\pi\xi$.

\begin{proposition}[Propriétés élémentaires]\label{appA:prop:fourierprop}
Soit $f$ intégrable, $a\in\R^d$, $\lambda>0$ et $M$ une matrice inversible.
\begin{enumerate}
\item $\hat f$ est continue et bornée par $\|f\|_1$.
\item Translation : $\widehat{f(\cdot-a)}(\xi)=e^{-2\pi i\,a\cdot\xi}\hat f(\xi)$.
\item Changement linéaire de variable : $\widehat{f\circ M}(\xi)=|\det M|^{-1}\hat f(M^{-\mathsf T}\xi)$.
\item Si $f$ est $\mathcal C^1$ et que $f$ et $\partial_jf$ sont intégrables, $\widehat{\partial_jf}(\xi)=2\pi
i\,\xi_j\hat f(\xi)$.
\end{enumerate}
\end{proposition}

\begin{humanproof}
(1) découle du théorème~\ref{appA:thm:cont}, avec $|f(x)e^{-2\pi i\xi\cdot x}|=|f(x)|$. (2) et (3) sont des changements
de variables $y=x-a$ et $y=Mx$ (jacobien $|\det M|$, et $\xi\cdot M^{-1}y=(M^{-\mathsf T}\xi)\cdot y$). (4) est une
intégration par parties en $x_j$ ; le terme de bord tend vers $0$ le long d'une suite de bornes, car $f$ est intégrable
et a une dérivée intégrable.
\end{humanproof}

\begin{example}[Gaussiennes]\label{appA:ex:gauss}
Pour $\alpha>0$, $\widehat{e^{-\pi\alpha|x|^2}}(\xi)=\alpha^{-d/2}e^{-\pi|\xi|^2/\alpha}$. La preuve en dimension~1
(une équation différentielle pour $\hat f$, obtenue par dérivation sous l'intégrale) est l'étape~1 de la démonstration
du théorème~\ref{ch1:thm:poisson} ; on passe à la dimension $d$ par Fubini, car l'exponentielle se factorise. En
particulier $e^{-\pi|x|^2}$ est sa propre transformée de Fourier. Pour une gaussienne de covariance quelconque,
le point~(3) de la proposition~\ref{appA:prop:fourierprop} donne la formule générale.
\end{example}

\begin{theorem}[Inversion \tierL]\label{appA:thm:inversion}
Si $f$ et $\hat f$ sont intégrables, alors $f(x)=\int\hat f(\xi)e^{2\pi i\,\xi\cdot x}d\xi$ en tout point où $f$ est
continue.
\end{theorem}

\subsection{Fonctions radiales du plan et fonction de Bessel \texorpdfstring{$J_0$}{J0}}\label{appA:ssec:hankel}

Les potentiels des chapitres~2 et~3 ne dépendent que de la distance $r=|x|$. Leur transformée de Fourier plane est
alors radiale aussi, et elle s'écrit avec la fonction de Bessel $J_0$, définie par
\begin{equation}\label{appA:eq:J0int}
J_0(z)=\frac1{2\pi}\int_0^{2\pi}e^{iz\cos\theta}\,d\theta=\frac1\pi\int_0^\pi\cos(z\cos\theta)\,d\theta .
\end{equation}

\begin{proposition}[Transformée de Hankel d'ordre $0$]\label{appA:prop:hankel}
Si $f(x)=F(|x|)$ est intégrable sur $\R^2$, alors $\hat f(\xi)=2\pi\int_0^\infty F(r)\,J_0(2\pi|\xi|r)\,r\,dr$.
\end{proposition}

\begin{humanproof}
En coordonnées polaires $x=r(\cos\theta,\sin\theta)$, et en choisissant les axes de sorte que
$\xi=(|\xi|,0)$ (l'intégrale est invariante par rotation), $\xi\cdot x=|\xi|r\cos\theta$. Par Fubini,
\[\hat f(\xi)=\int_0^\infty F(r)\Bigl(\int_0^{2\pi}e^{-2\pi i|\xi|r\cos\theta}d\theta\Bigr)r\,dr
=2\pi\int_0^\infty F(r)J_0(2\pi|\xi|r)\,r\,dr ,\]
car $J_0$ est paire (remplacer $\theta$ par $\theta+\pi$ dans~\eqref{appA:eq:J0int}).
\end{humanproof}

\begin{remark}[Contrôle sur la gaussienne]
Avec $F(r)=e^{-\pi ar^2}$, la proposition et l'exemple~\ref{appA:ex:gauss} donnent
$2\pi\int_0^\infty e^{-\pi ar^2}J_0(2\pi kr)\,r\,dr=a^{-1}e^{-\pi k^2/a}>0$ : c'est la formule employée au
chapitre~3 pour opposer la gaussienne, à transformée partout positive, au potentiel en épaulement des atomes de Rydberg
habillés, dont la transformée plane prend des valeurs négatives. Le calcul numérique de cette dernière est la
transformée de Hankel ci-dessus, évaluée par quadrature \tierX.
\end{remark}

\subsection{La transformée de Laplace}

\begin{definition}
Soit $\mu$ une mesure positive sur $[0,\infty[$ (par exemple $\mu(da)=w(a)\,da$ avec $w\ge0$). Sa \emph{transformée de
Laplace} est $\mathcal L\mu(t)=\int_0^\infty e^{-ta}\,\mu(da)$, pour les $t$ où l'intégrale converge.
\end{definition}

\begin{proposition}[Une transformée de Laplace est complètement monotone]\label{appA:prop:laplacecm}
Si $\mathcal L\mu(t)<\infty$ pour tout $t>0$, alors $\mathcal L\mu$ est $\mathcal C^\infty$ sur $]0,\infty[$ et
$(-1)^m(\mathcal L\mu)^{(m)}(t)=\int_0^\infty a^me^{-ta}\mu(da)\ge0$.
\end{proposition}

\begin{humanproof}
Sur $[t_0,\infty[$ avec $t_0>0$, on a $a^me^{-ta}\le a^me^{-t_0a}\le C_m\,e^{-t_0a/2}$, où
$C_m=\sup_{a\ge0}a^me^{-t_0a/2}<\infty$. Cette fonction est $\mu$-intégrable puisque $\mathcal L\mu(t_0/2)<\infty$.
Le théorème~\ref{appA:thm:deriv} (version mesure, même preuve) s'applique $m$ fois.
\end{humanproof}
C'est le sens « facile » du théorème de Bernstein (théorème~\ref{ch1:thm:bernstein}, proposition~\ref{ch1:prop:facile}) ;
la réciproque, toute fonction complètement monotone est une transformée de Laplace, est de niveau \tierL.

\begin{example}[Trois transformées de Laplace utilisées dans le livre]\label{appA:ex:laplace}
\begin{enumerate}
\item \emph{Puissances.} Pour $s>0$ et $t>0$, $t^{-s}=\frac1{\Gamma(s)}\int_0^\infty a^{s-1}e^{-ta}\,da$. En effet, le
changement de variable $u=ta$ donne $\int_0^\infty a^{s-1}e^{-ta}da=t^{-s}\int_0^\infty u^{s-1}e^{-u}du=t^{-s}\Gamma(s)$.
Avec $t=r^2$ et $s=1/2$, on obtient $1/r=\pi^{-1/2}\int_0^\infty a^{-1/2}e^{-ar^2}da$ : le potentiel de Coulomb est
un mélange de gaussiennes. Le changement $a=u^2$ donne la forme $1/r=\frac2{\sqrt\pi}\int_0^\infty e^{-r^2u^2}du$,
point de départ de la décomposition d'Ewald du cas U1 (chapitre~5).
\item \emph{Exciton dipolaire.} En retranchant la même formule écrite pour $r^2+d^2$,
$\frac1r-\frac1{\sqrt{r^2+d^2}}=\frac2{\sqrt\pi}\int_0^\infty e^{-r^2u^2}\bigl(1-e^{-d^2u^2}\bigr)du$, avec un poids
positif : c'est la formule du chapitre~5 (cas U1) au facteur $2$ près.
\item \emph{Subordination.} Pour $\rho>0$, $e^{-\rho\sqrt t}=\int_0^\infty e^{-ta}\,\frac{\rho}{2\sqrt\pi}\,a^{-3/2}
e^{-\rho^2/(4a)}\,da$ \tierL{} ; c'est la formule que le chapitre~2 emploie pour le potentiel de Yukawa et la
forme de Laplace du potentiel de Keldysh.
\end{enumerate}
\end{example}

%=====================================================================================================================
\section{La formule sommatoire de Poisson}\label{appA:sec:poisson}

Le chapitre~1 démontre la formule de Poisson pour les gaussiennes. Voici l'énoncé général dont elle est un cas
particulier, en dimension~$1$ d'abord, puis pour un réseau quelconque.

\begin{theorem}[Formule de Poisson en dimension 1]\label{appA:thm:poisson1}
Soit $f:\R\to\C$ continue. On suppose qu'il existe $C>0$ et $\delta>0$ tels que
\[|f(x)|\le C(1+|x|)^{-1-\delta}\quad\text{et}\quad|\hat f(\xi)|\le C(1+|\xi|)^{-1-\delta}\qquad(x,\xi\in\R).\]
Alors, pour tout $x\in\R$, $\displaystyle\sum_{n\in\Z}f(x+n)=\sum_{k\in\Z}\hat f(k)\,e^{2\pi ikx}$, les deux séries
convergeant absolument. En particulier $\sum_nf(n)=\sum_k\hat f(k)$.
\end{theorem}

\begin{humanproof}
\emph{Étape 1 : la périodisée.} Posons $F(x)=\sum_{n\in\Z}f(x+n)$. Sur $[0,1]$, $|f(x+n)|\le C(1+|n|-1)^{-1-\delta}$
pour $|n|\ge2$, terme général d'une série convergente ; la série converge donc normalement sur $[0,1]$, et par
périodicité sur $\R$. La fonction $F$ est continue (proposition~\ref{appA:prop:normale}) et $1$-périodique.

\emph{Étape 2 : ses coefficients de Fourier.} Pour $k\in\Z$, par le théorème~\ref{appA:thm:serieint} (la série des
intégrales des modules converge, car elle vaut $\int_\R|f|<\infty$),
\[c_k(F)=\int_0^1F(x)e^{-2\pi ikx}dx=\sum_{n\in\Z}\int_n^{n+1}f(y)e^{-2\pi ik(y-n)}dy=\int_\R f(y)e^{-2\pi iky}dy=\hat
f(k),\]
puisque $e^{2\pi ikn}=1$.

\emph{Étape 3 : la série de Fourier converge vers $F$.} L'hypothèse sur $\hat f$ donne $\sum_k|c_k(F)|<\infty$. La série
$G(x)=\sum_k\hat f(k)e^{2\pi ikx}$ converge donc normalement vers une fonction continue $1$-périodique, et
$c_k(G)=\hat f(k)$ (on intègre terme à terme). La fonction $H=F-G$ est continue, périodique, et tous ses coefficients
de Fourier sont nuls. Par le théorème de Fejér \tierL, les moyennes de Cesàro de la série de Fourier d'une fonction
continue périodique convergent uniformément vers elle ; ici elles sont toutes nulles, donc $H=0$.
\end{humanproof}

\begin{remark}[Pourquoi deux hypothèses de décroissance]
La décroissance de $f$ rend $F$ continue ; celle de $\hat f$ rend la série de Fourier absolument convergente. Sans la
seconde, la formule peut tomber en défaut : il existe des fonctions continues intégrables dont la série $\sum f(n)$
converge mais ne vaut pas $\sum\hat f(k)$. Les gaussiennes, et plus généralement les fonctions de Schwartz, vérifient
les deux hypothèses avec n'importe quel $\delta$.
\end{remark}

\begin{theorem}[Formule de Poisson pour un réseau]\label{appA:thm:poissonres}
Soit $\Lambda=M\Z^d$ un réseau de $\R^d$, de covolume $V=|\det M|$, et $\Lambda^*=M^{-\mathsf T}\Z^d$ son réseau dual.
Si $f$ est continue avec $|f(x)|\le C(1+|x|)^{-d-\delta}$ et $|\hat f(\xi)|\le C(1+|\xi|)^{-d-\delta}$, alors
\[\sum_{x\in\Lambda}f(x+y)=\frac1V\sum_{\xi\in\Lambda^*}\hat f(\xi)\,e^{2\pi i\,\xi\cdot y}\qquad(y\in\R^d).\]
\end{theorem}

\begin{humanproof}
On applique la version $d$-dimensionnelle de la démonstration précédente (le tore $\R^d/\Z^d$ à la place de
$\R/\Z$, et le théorème de Fejér sur le tore \tierL) à $g=f\circ M$. Les hypothèses de décroissance passent de $f$ à $g$,
car $|Mx|\ge\|M^{-1}\|^{-1}|x|$. On obtient $\sum_{n\in\Z^d}g(n+z)=\sum_{k\in\Z^d}\hat g(k)e^{2\pi ik\cdot z}$. Puis
$\hat g(k)=V^{-1}\hat f(M^{-\mathsf T}k)$ (proposition~\ref{appA:prop:fourierprop}, point~3). On pose enfin $y=Mz$ :
$k\cdot z=(M^{-\mathsf T}k)\cdot y$, et $M^{-\mathsf T}k$ parcourt $\Lambda^*$.
\end{humanproof}

Avec $f(x)=e^{-\pi\alpha|x|^2}$ et $y=0$, on retrouve exactement le théorème~\ref{ch1:thm:poisson}.

\subsection{Application : l'idée de la sommation d'Ewald}\label{appA:ssec:ewald}

Les sommes de réseau d'un potentiel à longue portée, comme $\sum_{x\in\Lambda\setminus\{0\}}|x|^{-3}$, convergent
lentement : l'erreur de troncature à la distance $R$ décroît comme $1/R$. La méthode d'Ewald coupe le potentiel en deux
grâce à une représentation en mélange de gaussiennes $V(r)=\int_0^\infty e^{-r^2u^2}w(u)\,du$, avec $w\ge0$
(exemple~\ref{appA:ex:laplace}). On choisit un paramètre $\eta>0$ et l'on écrit
\[V(r)=\underbrace{\int_\eta^\infty e^{-r^2u^2}w(u)\,du}_{\text{courte portée}}+
\underbrace{\int_0^\eta e^{-r^2u^2}w(u)\,du}_{\text{lisse, longue portée}} .\]
\begin{itemize}
\item La partie de courte portée décroît comme $e^{-\eta^2r^2}$ : sa somme sur le réseau converge très vite dans
l'espace direct.
\item La partie lisse est un mélange de gaussiennes larges. Le théorème~\ref{appA:thm:poissonres} transforme sa somme
en une somme sur le réseau dual de gaussiennes étroites $e^{-\pi^2|\xi|^2/u^2}$, qui converge aussi très vite.
\end{itemize}
Le résultat ne dépend pas de $\eta$ : c'est un contrôle naturel du code. Le cas U1 du chapitre~5 rapporte
justement que l'énergie d'Ewald de rusty-SUNDIALS ne varie pas avec ce paramètre de découpage à mieux que
$9{,}0\times10^{-15}$ près \tierX.

%=====================================================================================================================
\section{Fonctions spéciales}\label{appA:sec:speciales}

Nous ne démontrons ici que les propriétés élémentaires. Les autres sont citées du \emph{Digital Library of
Mathematical Functions} du NIST~\cite{app:dlmf} (chapitres~5 pour $\Gamma$, 7 pour $\operatorname{erf}$, 10 pour les
fonctions de Bessel, 11 pour celles de Struve) \tierL.

\subsection{La fonction \texorpdfstring{$\Gamma$}{Gamma}}
Pour $s>0$, $\Gamma(s)=\int_0^\infty u^{s-1}e^{-u}du$. Une intégration par parties donne $\Gamma(s+1)=s\Gamma(s)$, donc
$\Gamma(n+1)=n!$. Le changement $u=\pi x^2$ donne $\Gamma(1/2)=\int_\R e^{-\pi x^2}dx\cdot\sqrt\pi=\sqrt\pi$, grâce à
l'intégrale de Gauss.

\subsection{Fonction d'erreur et fonction d'erreur complémentaire}
\begin{definition}
$\operatorname{erf}(x)=\frac2{\sqrt\pi}\int_0^xe^{-s^2}ds$, $\operatorname{erfc}(x)=1-\operatorname{erf}(x)=
\frac2{\sqrt\pi}\int_x^\infty e^{-s^2}ds$, et $\operatorname{erfcx}(x)=e^{x^2}\operatorname{erfc}(x)$.
\end{definition}

\begin{proposition}\label{appA:prop:erfc}
Pour $x>0$ : $\displaystyle\frac{2}{\sqrt\pi}\,\frac{e^{-x^2}}{x+\sqrt{x^2+2}}\le\operatorname{erfc}(x)
\le\frac{e^{-x^2}}{x\sqrt\pi}$. En particulier $\operatorname{erfcx}(x)\le\frac1{x\sqrt\pi}$ et
$\operatorname{erfc}(x)\sim\frac{e^{-x^2}}{x\sqrt\pi}$ quand $x\to\infty$.
\end{proposition}

\begin{humanproof}
\emph{Majoration.} Pour $s\ge x>0$, $1\le s/x$, donc $\int_x^\infty e^{-s^2}ds\le\frac1x\int_x^\infty se^{-s^2}ds=
\frac{e^{-x^2}}{2x}$. On multiplie par $2/\sqrt\pi$.

\emph{Minoration.} On l'admet \tierL{} (elle se démontre en étudiant la différence des deux membres, qui tend vers $0$ à
l'infini et dont on contrôle le signe de la dérivée) ; elle n'est utilisée que pour l'équivalent, que la majoration et
une intégration par parties suffisent aussi à établir (exercice~\ref{appA:ex:erfc}).
\end{humanproof}
Le chapitre~2 utilise les bornes $0<\operatorname{erfcx}(x)\le1$ pour $x\ge0$ et $\operatorname{erfcx}(x)\le1/(\sqrt\pi
x)$ pour $x>0$, démontrées là par la représentation $\operatorname{erfcx}(x)=\frac2{\sqrt\pi}\int_0^\infty
e^{-(y^2+2xy)}dy$. Dans les sommes de réseau certifiées, $\operatorname{erfc}$ apparaît aussi dans les bornes de queue
des gaussiennes : $\sum_{|n|>R}e^{-\pi n^2}$ se compare à une intégrale de gaussienne, c'est-à-dire à
$\operatorname{erfc}(\sqrt\pi R)$.

\subsection{Les fonctions de Bessel \texorpdfstring{$J_0$}{J0}, \texorpdfstring{$Y_0$}{Y0} et \texorpdfstring{$K_0$}{K0}}

\begin{proposition}[$J_0$]\label{appA:prop:J0}
La fonction $J_0$ définie par~\eqref{appA:eq:J0int} vérifie :
\begin{enumerate}
\item $J_0(z)=\sum_{m\ge0}\frac{(-1)^m}{(m!)^2}\bigl(\frac z2\bigr)^{2m}$ pour tout $z\in\R$ ;
\item $|J_0(z)|\le1$ et $J_0(0)=1$ ;
\item $J_0$ est solution de l'équation de Bessel d'ordre $0$ : $zy''+y'+zy=0$ ;
\item \tierL{} $J_0(z)=\sqrt{2/(\pi z)}\,\bigl(\cos(z-\pi/4)+O(1/z)\bigr)$ quand $z\to+\infty$.
\end{enumerate}
\end{proposition}

\begin{humanproof}
(1) On développe $\cos(z\cos\theta)=\sum_m(-1)^m z^{2m}\cos^{2m}\theta/(2m)!$ ; la série converge normalement en
$\theta$, et l'on intègre terme à terme. Les intégrales de Wallis donnent
$\frac1\pi\int_0^\pi\cos^{2m}\theta\,d\theta=\binom{2m}m4^{-m}$, d'où le coefficient
$\frac{1}{(2m)!}\binom{2m}{m}4^{-m}=\frac{1}{(m!)^2 4^m}$.
(2) Le module de l'intégrande de~\eqref{appA:eq:J0int} vaut $1$.
(3) Par dérivation sous l'intégrale, $J_0'(z)=-\frac1\pi\int_0^\pi\cos\theta\sin(z\cos\theta)d\theta$ et
$J_0''(z)=-\frac1\pi\int_0^\pi\cos^2\theta\cos(z\cos\theta)d\theta$. Donc
$zJ_0''+zJ_0=\frac z\pi\int_0^\pi\sin^2\theta\cos(z\cos\theta)d\theta$. Une intégration par parties, avec
$\frac{d}{d\theta}\sin(z\cos\theta)=-z\sin\theta\cos(z\cos\theta)$, transforme cette intégrale en
$\frac1\pi\int_0^\pi\cos\theta\sin(z\cos\theta)d\theta=-J_0'(z)$.
\end{humanproof}

L'équation de Bessel a une seconde solution, singulière en $0$ : la fonction de Neumann $Y_0$, avec
$Y_0(z)\sim\frac2\pi\ln z$ quand $z\to0^+$ \tierL. La \emph{fonction de Bessel modifiée} $K_0$ est définie pour $x>0$ par
\[K_0(x)=\int_0^\infty e^{-x\cosh t}\,dt .\]
Elle est positive, décroissante, se comporte comme $-\ln x$ en $0$ et comme $\sqrt{\pi/(2x)}\,e^{-x}$ à l'infini \tierL.
C'est la fonction de Green de l'opérateur $-\Delta+\kappa^2$ dans le plan : $(-\Delta+\kappa^2)\frac1{2\pi}K_0(\kappa|x|)=\delta_0$
\tierL. Les chapitres~2 et~5 notent que ce potentiel de Debye--Hückel bidimensionnel n'est \emph{pas} traité dans
ce livre, qui étudie le Yukawa tridimensionnel restreint au plan, $e^{-\kappa r}/r$.

\subsection{La fonction de Struve \texorpdfstring{$\mathbf H_0$}{H0} et le potentiel de Rytova--Keldysh}

\begin{definition}
La fonction de Struve d'ordre $0$ est $\mathbf H_0(z)=\frac2\pi\sum_{m\ge0}\frac{(-1)^mz^{2m+1}}{\bigl(1\cdot3\cdot5\cdots(2m+1)\bigr)^2}$
pour $z\in\R$.
\end{definition}
La série converge pour tout $z$ (critère de d'Alembert). On a aussi la représentation intégrale
$\mathbf H_0(z)=\frac2\pi\int_0^{\pi/2}\sin(z\cos\theta)\,d\theta$ \tierL, et $\mathbf H_0(z)-Y_0(z)\sim\frac2{\pi z}$
quand $z\to+\infty$ \tierL. Mathlib, à la version épinglée, ne connaît ni $\mathbf H_0$ ni $Y_0$ : c'est pourquoi le
chapitre~2 formalise le potentiel de Keldysh sous une forme de Laplace, et relie cette forme à la combinaison
$\mathbf H_0-Y_0$ par un argument de niveau \tierC{} qui s'appuie sur la représentation
$\mathbf K_0(z)=\mathbf H_0(z)-Y_0(z)=\frac2\pi\int_0^\infty e^{-zu}(1+u^2)^{-1/2}du$ pour $\operatorname{Re}z>0$
(DLMF~11.5.2, citée au chapitre~2~\cite{ch2dlmf}). Le comportement $\sim 2/(\pi z)$ est la loi de Coulomb $1/r$
retrouvée aux grandes distances.

%=====================================================================================================================
\section{Réseaux du plan et domaine fondamental de \texorpdfstring{$SL_2(\Z)$}{SL2(Z)}}\label{appA:sec:sl2}

\subsection{Bases d'un réseau}
Un réseau du plan est $\Lambda=\Z u\oplus\Z v$ avec $(u,v)$ base de $\R^2$ (définition~\ref{ch1:def:reseau}). Deux bases
$(u,v)$ et $(u',v')$ engendrent le même réseau si et seulement si $\binom{u'}{v'}=P\binom uv$ avec $P\in GL_2(\Z)$,
c'est-à-dire $P$ à coefficients entiers et $\det P=\pm1$.

\begin{humanproof}
Si $\Z u'\oplus\Z v'=\Z u\oplus\Z v$, chaque base s'écrit dans l'autre avec des coefficients entiers : $P$ et $P'$
entières avec $PP'=I$, donc $\det P\det P'=1$ dans $\Z$, d'où $\det P=\pm1$. Réciproquement, si $P\in GL_2(\Z)$, alors
$P^{-1}=\pm\,\mathrm{com}(P)^{\mathsf T}$ est entière, et les deux réseaux sont inclus l'un dans l'autre.
\end{humanproof}
Le covolume $|\det(u,v)|$ ne dépend donc pas de la base.

\subsection{Le paramètre de forme \texorpdfstring{$\tau$}{tau}}
Identifions $\R^2$ à $\C$. Quitte à faire une similitude directe (une rotation et une homothétie, qui ne changent pas la
« forme »), on peut supposer $u=1$. Quitte à remplacer $v$ par $-v$, on peut supposer $\operatorname{Im}v>0$. Le réseau
est alors $\Z+\Z\tau$ avec $\tau$ dans le demi-plan de Poincaré $\mathbb H=\{\operatorname{Im}\tau>0\}$. Un changement
de base direct $\binom{u'}{v'}=\begin{pmatrix}d&c\\b&a\end{pmatrix}\binom uv$ avec $ad-bc=1$ remplace $\tau$ par
\[\gamma\cdot\tau=\frac{a\tau+b}{c\tau+d},\qquad\gamma=\begin{pmatrix}a&b\\c&d\end{pmatrix}\in SL_2(\Z),\]
après renormalisation de $u'$ à $1$. Les formes de réseaux à similitude près correspondent donc aux orbites de
$SL_2(\Z)$ dans $\mathbb H$. Le calcul clé est
\begin{equation}\label{appA:eq:im}
\operatorname{Im}(\gamma\cdot\tau)=\frac{\operatorname{Im}\tau}{|c\tau+d|^2},
\end{equation}
obtenu en multipliant numérateur et dénominateur par $\overline{c\tau+d}$ et en utilisant $ad-bc=1$.

\subsection{Le domaine fondamental}
\begin{theorem}\label{appA:thm:domaine}
Soit $\mathcal F=\{\tau\in\mathbb H:|\operatorname{Re}\tau|\le\tfrac12,\ |\tau|\ge1\}$.
\begin{enumerate}
\item Tout $\tau\in\mathbb H$ a un point de son orbite sous $SL_2(\Z)$ dans $\mathcal F$.
\item Si deux points distincts de $\mathcal F$ sont dans la même orbite, ils sont sur le bord de $\mathcal F$
(symétriques par rapport à l'axe imaginaire, sur les droites $\operatorname{Re}\tau=\pm\frac12$ ou sur le cercle
$|\tau|=1$).
\end{enumerate}
\end{theorem}

\begin{humanproof}[Démonstration humaine du point 1]
Notons $T:\tau\mapsto\tau+1$ et $S:\tau\mapsto-1/\tau$ ; ils correspondent à $\bigl(\begin{smallmatrix}1&1\\0&1\end{smallmatrix}\bigr)$
et $\bigl(\begin{smallmatrix}0&-1\\1&0\end{smallmatrix}\bigr)$. Pour $\tau$ fixé, les nombres $c\tau+d$, $(c,d)\in\Z^2\setminus\{0\}$,
forment un réseau de $\C$ ; il n'y en a qu'un nombre fini de module $\le1$. D'après~\eqref{appA:eq:im},
$\operatorname{Im}(\gamma\cdot\tau)$ ne prend donc qu'un nombre fini de valeurs $\ge\operatorname{Im}\tau$, et il existe
$\gamma_0$ telle que $\tau'=\gamma_0\cdot\tau$ ait une partie imaginaire maximale dans l'orbite. Une puissance de $T$
amène $\operatorname{Re}\tau'$ dans $[-\frac12,\frac12]$ sans changer sa partie imaginaire. Si l'on avait $|\tau'|<1$,
alors $\operatorname{Im}(S\tau')=\operatorname{Im}\tau'/|\tau'|^2>\operatorname{Im}\tau'$, ce qui contredirait la
maximalité. Donc $\tau'\in\mathcal F$.
\end{humanproof}

Le point~2, dont la preuve est un examen de cas à partir de~\eqref{appA:eq:im}, est de niveau \tierL~\cite{app:serre}.
Le domaine $\mathcal F$ contient deux points remarquables :
\begin{itemize}
\item $\tau=i$, le réseau carré. Son stabilisateur dans $SL_2(\Z)$ est engendré par $S$, d'ordre $4$ dans $SL_2(\Z)$
(et $2$ dans le quotient par $\pm I$) : $S$ agit au voisinage de $i$ comme une rotation d'angle $\pi$.
\item $\tau=\rho=e^{i\pi/3}=\frac12+i\frac{\sqrt3}2$, le réseau triangulaire. Il est fixé par
$\gamma(\tau)=(\tau-1)/\tau$, d'ordre $3$ modulo $\pm I$, qui agit au voisinage de $\rho$ comme une rotation d'angle
$\pm120^\circ$.
\end{itemize}
C'est l'argument du §\ref{ch3:sec:isotropie} : une fonction de la forme du réseau, invariante par $SL_2(\Z)$ et de
classe $\mathcal C^2$, a un gradient nul en $\rho$ et une hessienne scalaire, parce que la seule matrice symétrique qui
commute avec une rotation de $120^\circ$ est un multiple de l'identité. En $i$, la rotation d'angle $\pi$ ne force que
le gradient à s'annuler : la hessienne peut y être anisotrope.

\subsection{Réduction de Gauss d'une base}
Le point~1 du théorème est aussi un algorithme. Partant d'une base $(u,v)$ :
\begin{enumerate}
\item si $|v|<|u|$, échanger $u$ et $v$ ;
\item remplacer $v$ par $v-ku$, où $k$ est l'entier le plus proche de $\operatorname{Re}(v\bar u)/|u|^2$ ;
\item recommencer tant que $|v|<|u|$.
\end{enumerate}
L'algorithme termine, car $|u|$ décroît strictement à chaque échange et ne prend qu'un nombre fini de valeurs (ce sont
des normes de vecteurs du réseau inférieures à la norme initiale). À la fin, $|u|\le|v|\le|v\pm u|$ : $u$ est un plus
court vecteur non nul et $\tau=v/u$ (au signe près) est dans $\mathcal F$. Les sommes de réseau tronquées des chapitres~3
et~5 utilisent une base ainsi réduite, pour laquelle le disque de rayon $R$ ne contient que $O(R^2)$ points de la forme
$mu+nv$ avec $|m|,|n|\le CR$.

%=====================================================================================================================
\section{Exercices}\label{appA:sec:exos}

\begin{exercise}\label{appA:ex:moments}
Démontrer par récurrence sur $m$ la formule des moments gaussiens de l'exemple~\ref{appA:ex:moments}. En déduire
$\int_\R x^{2m}e^{-x^2}dx=\Gamma(m+\tfrac12)$.
\end{exercise}

\begin{exercise}\label{appA:ex:erfc}
Montrer, par une intégration par parties, que $\operatorname{erfc}(x)=\frac{e^{-x^2}}{x\sqrt\pi}\bigl(1-\frac1{2x^2}
+O(x^{-4})\bigr)$ quand $x\to\infty$. Combien de termes de $\sum_{n\ge1}e^{-\pi n^2}$ faut-il pour une erreur inférieure
à $10^{-30}$ ?
\end{exercise}

\begin{exercise}
Appliquer le théorème~\ref{appA:thm:poisson1} à $f(x)=e^{-2\pi a|x|}$, $a>0$, dont la transformée de Fourier est
$\hat f(\xi)=\frac{a}{\pi(a^2+\xi^2)}$ (à vérifier). En déduire $\sum_{k\in\Z}\frac1{a^2+k^2}=\frac\pi a\coth(\pi a)$,
puis $\sum_{k\ge1}k^{-2}=\pi^2/6$ en faisant tendre $a$ vers $0$.
\end{exercise}

\begin{exercise}
Vérifier numériquement la proposition~\ref{appA:prop:hankel} pour $F(r)=e^{-\pi r^2}$ en évaluant $J_0$ par sa série
entière (avec \texttt{mpmath}, par exemple). Pourquoi la série entière de $J_0$ est-elle un mauvais moyen de calculer
$J_0(50)$ en virgule flottante ? (Penser à l'annulation de termes de l'ordre de $e^{50}$.)
\end{exercise}

\begin{exercise}
Réduire par l'algorithme de Gauss la base $u=(1,0)$, $v=(7{,}5;\,0{,}2)$, et donner le paramètre $\tau\in\mathcal F$
obtenu. Calculer le stabilisateur de $\rho$ dans $SL_2(\Z)$ et vérifier qu'il a six éléments.
\end{exercise}
.
% Etiquettes prefixees appA: ; references bibliographiques prefixees app: (refs_app.tex).
\chapter{Rappels d'analyse}\label{appA:chap}

Cette annexe rassemble les outils d'analyse que les chapitres utilisent sans toujours les redémontrer :
\begin{itemize}
\item l'interversion d'une limite et d'une intégrale ou d'une série (§\ref{appA:sec:param}) ;
\item les transformées de Fourier et de Laplace (§\ref{appA:sec:fourier}) ;
\item la formule sommatoire de Poisson, dans le cas général (§\ref{appA:sec:poisson}) ;
\item les fonctions spéciales rencontrées : $\Gamma$, $\operatorname{erfc}$, $J_0$, $K_0$, $Y_0$ et la fonction de
Struve $\mathbf H_0$ (§\ref{appA:sec:speciales}) ;
\item les réseaux du plan et le domaine fondamental de $SL_2(\Z)$ (§\ref{appA:sec:sl2}).
\end{itemize}
Le niveau visé est celui de la deuxième année de classe préparatoire. L'intégrale est celle de Lebesgue, mais il suffit
de savoir appliquer les deux théorèmes de convergence du programme.

\begin{warnbox}[: niveaux des énoncés de cette annexe]
Les démonstrations écrites ici sont des démonstrations humaines \tierC{} : aucune n'est vérifiée par Lean dans ce livre.
Les énoncés admis (convergence dominée, théorème de Fejér, propriétés des fonctions de Bessel et de Struve tirées du
DLMF~\cite{app:dlmf}) sont de niveau \tierL. La Mathlib épinglée contient des versions formelles de plusieurs d'entre
eux : la convergence dominée (\lean{MeasureTheory.tendsto_integral_of_dominated_convergence}, fichier
\texttt{MeasureTheory/Integral/DominatedConvergence.lean}) et la formule de Poisson (fichiers
\texttt{Analysis/Fourier/PoissonSummation.lean} et \texttt{Analysis/SpecialFunctions/Gaussian/PoissonSummation.lean}).
Nous avons seulement constaté que ces fichiers existent ; nous n'avons ni comparé leurs énoncés à ceux de cette annexe,
ni relié ces énoncés aux démonstrations ci-dessous.
\end{warnbox}

%=====================================================================================================================
\section{Séries et intégrales dépendant d'un paramètre}\label{appA:sec:param}

\subsection{Le théorème de convergence dominée}

\begin{theorem}[Convergence dominée \tierL]\label{appA:thm:cvd}
Soit $I$ un intervalle de $\R$ et $(f_n)$ une suite de fonctions continues par morceaux sur $I$. On suppose que :
\begin{enumerate}
\item $(f_n)$ converge simplement sur $I$ vers une fonction $f$ continue par morceaux ;
\item il existe $\varphi:I\to\R_+$ intégrable sur $I$ telle que $|f_n(x)|\le\varphi(x)$ pour tout $n$ et tout $x$
(hypothèse de domination).
\end{enumerate}
Alors les $f_n$ et $f$ sont intégrables sur $I$, et $\int_If_n\to\int_If$.
\end{theorem}

Le point essentiel est que la fonction dominante $\varphi$ ne dépend pas de $n$. Sans elle, la conclusion peut échouer :
$f_n=n\,\mathbf 1_{]0,1/n]}$ tend simplement vers $0$ sur $]0,1]$, alors que $\int f_n=1$ pour tout $n$.

\begin{theorem}[Interversion série--intégrale \tierL]\label{appA:thm:serieint}
Soit $(u_n)$ une suite de fonctions intégrables sur $I$ telle que $\sum_n u_n$ converge simplement vers une fonction
continue par morceaux et $\sum_n\int_I|u_n|<\infty$. Alors $\sum_nu_n$ est intégrable et
$\int_I\sum_nu_n=\sum_n\int_Iu_n$.
\end{theorem}

\subsection{Continuité et dérivation sous l'intégrale}

\begin{theorem}[Continuité sous l'intégrale]\label{appA:thm:cont}
Soit $A$ une partie de $\R$ et $f:A\times I\to\C$ telle que $f(\cdot,x)$ est continue pour tout $x$, $f(a,\cdot)$ est
continue par morceaux pour tout $a$, et $|f(a,x)|\le\varphi(x)$ avec $\varphi$ intégrable. Alors
$F(a)=\int_If(a,x)\,dx$ est continue sur $A$.
\end{theorem}

\begin{humanproof}
Soit $a_n\to a$ dans $A$. Les fonctions $x\mapsto f(a_n,x)$ convergent simplement vers $f(a,\cdot)$ et sont dominées par
$\varphi$. Le théorème~\ref{appA:thm:cvd} donne $F(a_n)\to F(a)$, et la caractérisation séquentielle conclut.
\end{humanproof}

\begin{theorem}[Dérivation sous l'intégrale]\label{appA:thm:deriv}
Soit $J$ un intervalle et $f:J\times I\to\C$. On suppose que $f(a,\cdot)$ est intégrable pour tout $a\in J$, que
$\partial f/\partial a$ existe, est continue en $a$ et continue par morceaux en $x$, et qu'il existe $\psi$ intégrable
avec $|\partial_af(a,x)|\le\psi(x)$ pour tous $a,x$. Alors $F(a)=\int_If(a,x)\,dx$ est de classe $\mathcal C^1$ sur $J$,
et $F'(a)=\int_I\partial_af(a,x)\,dx$.
\end{theorem}

\begin{humanproof}
Fixons $a\in J$ et une suite $h_n\to0$, $h_n\ne0$, avec $a+h_n\in J$. Posons
$g_n(x)=\bigl(f(a+h_n,x)-f(a,x)\bigr)/h_n$. Par définition de la dérivée partielle, $g_n(x)\to\partial_af(a,x)$ pour
tout $x$. Par l'inégalité des accroissements finis appliquée à $t\mapsto f(t,x)$ entre $a$ et $a+h_n$,
$|g_n(x)|\le\sup_t|\partial_af(t,x)|\le\psi(x)$. Le théorème~\ref{appA:thm:cvd} donne
$\bigl(F(a+h_n)-F(a)\bigr)/h_n=\int g_n\to\int\partial_af(a,x)\,dx$. La continuité de $F'$ découle du
théorème~\ref{appA:thm:cont} appliqué à $\partial_af$.
\end{humanproof}

\begin{remark}[Domination locale]
Il suffit en pratique que la domination ait lieu sur tout segment $[\alpha,\beta]\subset J$ : la dérivabilité est une
propriété locale. C'est ainsi qu'on traite les transformées de Laplace ci-dessous, où la domination n'est uniforme que
pour $t\ge\alpha>0$.
\end{remark}

\begin{example}[Moments gaussiens]\label{appA:ex:moments}
Pour $\alpha>0$, $G(\alpha)=\int_\R e^{-\pi\alpha x^2}dx=\alpha^{-1/2}$ (intégrale de Gauss). Sur $[\alpha_0,\infty[$ avec
$\alpha_0>0$, $|\partial_\alpha e^{-\pi\alpha x^2}|=\pi x^2e^{-\pi\alpha x^2}\le\pi x^2e^{-\pi\alpha_0x^2}$, qui est
intégrable. Le théorème~\ref{appA:thm:deriv} donne $\int\pi x^2e^{-\pi\alpha x^2}dx=-G'(\alpha)=\tfrac12\alpha^{-3/2}$.
En dérivant $m$ fois, $\int(\pi x^2)^me^{-\pi\alpha x^2}dx=\frac{(2m)!}{4^mm!}\,\alpha^{-1/2-m}$ (exercice~\ref{appA:ex:moments}).
C'est la même technique qui, au chapitre~1, calcule la transformée de Fourier d'une gaussienne
(théorème~\ref{ch1:thm:poisson}, étape~1).
\end{example}

\subsection{Séries de fonctions}

\begin{proposition}[Convergence normale \tierL]\label{appA:prop:normale}
Soit $(u_n)$ des fonctions continues sur une partie $A$ de $\R^d$ avec $\sum_n\sup_A|u_n|<\infty$. Alors
$\sum_nu_n$ converge uniformément sur $A$ et sa somme est continue. Si de plus les $u_n$ sont $\mathcal C^1$ sur un
ouvert et $\sum_n\sup|\nabla u_n|<\infty$, la somme est $\mathcal C^1$ et se dérive terme à terme.
\end{proposition}

\begin{example}[Séries thêta]
Pour un réseau $\Lambda$ du plan et $\alpha\ge\alpha_0>0$, on a $|x|^2\ge c\,(m^2+n^2)$ pour $x=mu+nv$ (chapitre~1),
donc $e^{-\pi\alpha|x|^2}\le e^{-\pi\alpha_0c(m^2+n^2)}$, terme général d'une série convergente. La série
$\Theta_\Lambda(\alpha)$ converge donc normalement sur $[\alpha_0,\infty[$, et elle y est $\mathcal C^\infty$ en
$\alpha$ : on dérive terme à terme autant qu'on veut. Les sommes de réseau $S(\alpha)=\sum|P|^2e^{-\pi\alpha|P|^2}$ du
chapitre~3 sont précisément $-\Theta_\Lambda'(\alpha)/\pi$.
\end{example}

%=====================================================================================================================
\section{Transformées de Fourier et de Laplace}\label{appA:sec:fourier}

\subsection{La transformée de Fourier}

\begin{definition}
Pour $f:\R^d\to\C$ intégrable, la \emph{transformée de Fourier} de $f$ est
\[\hat f(\xi)=\int_{\R^d}f(x)\,e^{-2\pi i\,\xi\cdot x}\,dx,\qquad\xi\in\R^d .\]
\end{definition}
C'est la convention du chapitre~1 : le facteur $2\pi$ est dans l'exponentielle, et la formule d'inversion n'a alors
aucun facteur devant l'intégrale. D'autres livres, en physique surtout, écrivent $\int f(x)e^{-ik\cdot x}dx$ ; on passe
de l'une à l'autre par $k=2\pi\xi$.

\begin{proposition}[Propriétés élémentaires]\label{appA:prop:fourierprop}
Soit $f$ intégrable, $a\in\R^d$, $\lambda>0$ et $M$ une matrice inversible.
\begin{enumerate}
\item $\hat f$ est continue et bornée par $\|f\|_1$.
\item Translation : $\widehat{f(\cdot-a)}(\xi)=e^{-2\pi i\,a\cdot\xi}\hat f(\xi)$.
\item Changement linéaire de variable : $\widehat{f\circ M}(\xi)=|\det M|^{-1}\hat f(M^{-\mathsf T}\xi)$.
\item Si $f$ est $\mathcal C^1$ et que $f$ et $\partial_jf$ sont intégrables, $\widehat{\partial_jf}(\xi)=2\pi
i\,\xi_j\hat f(\xi)$.
\end{enumerate}
\end{proposition}

\begin{humanproof}
(1) découle du théorème~\ref{appA:thm:cont}, avec $|f(x)e^{-2\pi i\xi\cdot x}|=|f(x)|$. (2) et (3) sont des changements
de variables $y=x-a$ et $y=Mx$ (jacobien $|\det M|$, et $\xi\cdot M^{-1}y=(M^{-\mathsf T}\xi)\cdot y$). (4) est une
intégration par parties en $x_j$ ; le terme de bord tend vers $0$ le long d'une suite de bornes, car $f$ est intégrable
et a une dérivée intégrable.
\end{humanproof}

\begin{example}[Gaussiennes]\label{appA:ex:gauss}
Pour $\alpha>0$, $\widehat{e^{-\pi\alpha|x|^2}}(\xi)=\alpha^{-d/2}e^{-\pi|\xi|^2/\alpha}$. La preuve en dimension~1
(une équation différentielle pour $\hat f$, obtenue par dérivation sous l'intégrale) est l'étape~1 de la démonstration
du théorème~\ref{ch1:thm:poisson} ; on passe à la dimension $d$ par Fubini, car l'exponentielle se factorise. En
particulier $e^{-\pi|x|^2}$ est sa propre transformée de Fourier. Pour une gaussienne de covariance quelconque,
le point~(3) de la proposition~\ref{appA:prop:fourierprop} donne la formule générale.
\end{example}

\begin{theorem}[Inversion \tierL]\label{appA:thm:inversion}
Si $f$ et $\hat f$ sont intégrables, alors $f(x)=\int\hat f(\xi)e^{2\pi i\,\xi\cdot x}d\xi$ en tout point où $f$ est
continue.
\end{theorem}

\subsection{Fonctions radiales du plan et fonction de Bessel \texorpdfstring{$J_0$}{J0}}\label{appA:ssec:hankel}

Les potentiels des chapitres~2 et~3 ne dépendent que de la distance $r=|x|$. Leur transformée de Fourier plane est
alors radiale aussi, et elle s'écrit avec la fonction de Bessel $J_0$, définie par
\begin{equation}\label{appA:eq:J0int}
J_0(z)=\frac1{2\pi}\int_0^{2\pi}e^{iz\cos\theta}\,d\theta=\frac1\pi\int_0^\pi\cos(z\cos\theta)\,d\theta .
\end{equation}

\begin{proposition}[Transformée de Hankel d'ordre $0$]\label{appA:prop:hankel}
Si $f(x)=F(|x|)$ est intégrable sur $\R^2$, alors $\hat f(\xi)=2\pi\int_0^\infty F(r)\,J_0(2\pi|\xi|r)\,r\,dr$.
\end{proposition}

\begin{humanproof}
En coordonnées polaires $x=r(\cos\theta,\sin\theta)$, et en choisissant les axes de sorte que
$\xi=(|\xi|,0)$ (l'intégrale est invariante par rotation), $\xi\cdot x=|\xi|r\cos\theta$. Par Fubini,
\[\hat f(\xi)=\int_0^\infty F(r)\Bigl(\int_0^{2\pi}e^{-2\pi i|\xi|r\cos\theta}d\theta\Bigr)r\,dr
=2\pi\int_0^\infty F(r)J_0(2\pi|\xi|r)\,r\,dr ,\]
car $J_0$ est paire (remplacer $\theta$ par $\theta+\pi$ dans~\eqref{appA:eq:J0int}).
\end{humanproof}

\begin{remark}[Contrôle sur la gaussienne]
Avec $F(r)=e^{-\pi ar^2}$, la proposition et l'exemple~\ref{appA:ex:gauss} donnent
$2\pi\int_0^\infty e^{-\pi ar^2}J_0(2\pi kr)\,r\,dr=a^{-1}e^{-\pi k^2/a}>0$ : c'est la formule employée au
chapitre~3 pour opposer la gaussienne, à transformée partout positive, au potentiel en épaulement des atomes de Rydberg
habillés, dont la transformée plane prend des valeurs négatives. Le calcul numérique de cette dernière est la
transformée de Hankel ci-dessus, évaluée par quadrature \tierX.
\end{remark}

\subsection{La transformée de Laplace}

\begin{definition}
Soit $\mu$ une mesure positive sur $[0,\infty[$ (par exemple $\mu(da)=w(a)\,da$ avec $w\ge0$). Sa \emph{transformée de
Laplace} est $\mathcal L\mu(t)=\int_0^\infty e^{-ta}\,\mu(da)$, pour les $t$ où l'intégrale converge.
\end{definition}

\begin{proposition}[Une transformée de Laplace est complètement monotone]\label{appA:prop:laplacecm}
Si $\mathcal L\mu(t)<\infty$ pour tout $t>0$, alors $\mathcal L\mu$ est $\mathcal C^\infty$ sur $]0,\infty[$ et
$(-1)^m(\mathcal L\mu)^{(m)}(t)=\int_0^\infty a^me^{-ta}\mu(da)\ge0$.
\end{proposition}

\begin{humanproof}
Sur $[t_0,\infty[$ avec $t_0>0$, on a $a^me^{-ta}\le a^me^{-t_0a}\le C_m\,e^{-t_0a/2}$, où
$C_m=\sup_{a\ge0}a^me^{-t_0a/2}<\infty$. Cette fonction est $\mu$-intégrable puisque $\mathcal L\mu(t_0/2)<\infty$.
Le théorème~\ref{appA:thm:deriv} (version mesure, même preuve) s'applique $m$ fois.
\end{humanproof}
C'est le sens « facile » du théorème de Bernstein (théorème~\ref{ch1:thm:bernstein}, proposition~\ref{ch1:prop:facile}) ;
la réciproque, toute fonction complètement monotone est une transformée de Laplace, est de niveau \tierL.

\begin{example}[Trois transformées de Laplace utilisées dans le livre]\label{appA:ex:laplace}
\begin{enumerate}
\item \emph{Puissances.} Pour $s>0$ et $t>0$, $t^{-s}=\frac1{\Gamma(s)}\int_0^\infty a^{s-1}e^{-ta}\,da$. En effet, le
changement de variable $u=ta$ donne $\int_0^\infty a^{s-1}e^{-ta}da=t^{-s}\int_0^\infty u^{s-1}e^{-u}du=t^{-s}\Gamma(s)$.
Avec $t=r^2$ et $s=1/2$, on obtient $1/r=\pi^{-1/2}\int_0^\infty a^{-1/2}e^{-ar^2}da$ : le potentiel de Coulomb est
un mélange de gaussiennes. Le changement $a=u^2$ donne la forme $1/r=\frac2{\sqrt\pi}\int_0^\infty e^{-r^2u^2}du$,
point de départ de la décomposition d'Ewald du cas U1 (chapitre~5).
\item \emph{Exciton dipolaire.} En retranchant la même formule écrite pour $r^2+d^2$,
$\frac1r-\frac1{\sqrt{r^2+d^2}}=\frac2{\sqrt\pi}\int_0^\infty e^{-r^2u^2}\bigl(1-e^{-d^2u^2}\bigr)du$, avec un poids
positif : c'est la formule du chapitre~5 (cas U1) au facteur $2$ près.
\item \emph{Subordination.} Pour $\rho>0$, $e^{-\rho\sqrt t}=\int_0^\infty e^{-ta}\,\frac{\rho}{2\sqrt\pi}\,a^{-3/2}
e^{-\rho^2/(4a)}\,da$ \tierL{} ; c'est la formule que le chapitre~2 emploie pour le potentiel de Yukawa et la
forme de Laplace du potentiel de Keldysh.
\end{enumerate}
\end{example}

%=====================================================================================================================
\section{La formule sommatoire de Poisson}\label{appA:sec:poisson}

Le chapitre~1 démontre la formule de Poisson pour les gaussiennes. Voici l'énoncé général dont elle est un cas
particulier, en dimension~$1$ d'abord, puis pour un réseau quelconque.

\begin{theorem}[Formule de Poisson en dimension 1]\label{appA:thm:poisson1}
Soit $f:\R\to\C$ continue. On suppose qu'il existe $C>0$ et $\delta>0$ tels que
\[|f(x)|\le C(1+|x|)^{-1-\delta}\quad\text{et}\quad|\hat f(\xi)|\le C(1+|\xi|)^{-1-\delta}\qquad(x,\xi\in\R).\]
Alors, pour tout $x\in\R$, $\displaystyle\sum_{n\in\Z}f(x+n)=\sum_{k\in\Z}\hat f(k)\,e^{2\pi ikx}$, les deux séries
convergeant absolument. En particulier $\sum_nf(n)=\sum_k\hat f(k)$.
\end{theorem}

\begin{humanproof}
\emph{Étape 1 : la périodisée.} Posons $F(x)=\sum_{n\in\Z}f(x+n)$. Sur $[0,1]$, $|f(x+n)|\le C(1+|n|-1)^{-1-\delta}$
pour $|n|\ge2$, terme général d'une série convergente ; la série converge donc normalement sur $[0,1]$, et par
périodicité sur $\R$. La fonction $F$ est continue (proposition~\ref{appA:prop:normale}) et $1$-périodique.

\emph{Étape 2 : ses coefficients de Fourier.} Pour $k\in\Z$, par le théorème~\ref{appA:thm:serieint} (la série des
intégrales des modules converge, car elle vaut $\int_\R|f|<\infty$),
\[c_k(F)=\int_0^1F(x)e^{-2\pi ikx}dx=\sum_{n\in\Z}\int_n^{n+1}f(y)e^{-2\pi ik(y-n)}dy=\int_\R f(y)e^{-2\pi iky}dy=\hat
f(k),\]
puisque $e^{2\pi ikn}=1$.

\emph{Étape 3 : la série de Fourier converge vers $F$.} L'hypothèse sur $\hat f$ donne $\sum_k|c_k(F)|<\infty$. La série
$G(x)=\sum_k\hat f(k)e^{2\pi ikx}$ converge donc normalement vers une fonction continue $1$-périodique, et
$c_k(G)=\hat f(k)$ (on intègre terme à terme). La fonction $H=F-G$ est continue, périodique, et tous ses coefficients
de Fourier sont nuls. Par le théorème de Fejér \tierL, les moyennes de Cesàro de la série de Fourier d'une fonction
continue périodique convergent uniformément vers elle ; ici elles sont toutes nulles, donc $H=0$.
\end{humanproof}

\begin{remark}[Pourquoi deux hypothèses de décroissance]
La décroissance de $f$ rend $F$ continue ; celle de $\hat f$ rend la série de Fourier absolument convergente. Sans la
seconde, la formule peut tomber en défaut : il existe des fonctions continues intégrables dont la série $\sum f(n)$
converge mais ne vaut pas $\sum\hat f(k)$. Les gaussiennes, et plus généralement les fonctions de Schwartz, vérifient
les deux hypothèses avec n'importe quel $\delta$.
\end{remark}

\begin{theorem}[Formule de Poisson pour un réseau]\label{appA:thm:poissonres}
Soit $\Lambda=M\Z^d$ un réseau de $\R^d$, de covolume $V=|\det M|$, et $\Lambda^*=M^{-\mathsf T}\Z^d$ son réseau dual.
Si $f$ est continue avec $|f(x)|\le C(1+|x|)^{-d-\delta}$ et $|\hat f(\xi)|\le C(1+|\xi|)^{-d-\delta}$, alors
\[\sum_{x\in\Lambda}f(x+y)=\frac1V\sum_{\xi\in\Lambda^*}\hat f(\xi)\,e^{2\pi i\,\xi\cdot y}\qquad(y\in\R^d).\]
\end{theorem}

\begin{humanproof}
On applique la version $d$-dimensionnelle de la démonstration précédente (le tore $\R^d/\Z^d$ à la place de
$\R/\Z$, et le théorème de Fejér sur le tore \tierL) à $g=f\circ M$. Les hypothèses de décroissance passent de $f$ à $g$,
car $|Mx|\ge\|M^{-1}\|^{-1}|x|$. On obtient $\sum_{n\in\Z^d}g(n+z)=\sum_{k\in\Z^d}\hat g(k)e^{2\pi ik\cdot z}$. Puis
$\hat g(k)=V^{-1}\hat f(M^{-\mathsf T}k)$ (proposition~\ref{appA:prop:fourierprop}, point~3). On pose enfin $y=Mz$ :
$k\cdot z=(M^{-\mathsf T}k)\cdot y$, et $M^{-\mathsf T}k$ parcourt $\Lambda^*$.
\end{humanproof}

Avec $f(x)=e^{-\pi\alpha|x|^2}$ et $y=0$, on retrouve exactement le théorème~\ref{ch1:thm:poisson}.

\subsection{Application : l'idée de la sommation d'Ewald}\label{appA:ssec:ewald}

Les sommes de réseau d'un potentiel à longue portée, comme $\sum_{x\in\Lambda\setminus\{0\}}|x|^{-3}$, convergent
lentement : l'erreur de troncature à la distance $R$ décroît comme $1/R$. La méthode d'Ewald coupe le potentiel en deux
grâce à une représentation en mélange de gaussiennes $V(r)=\int_0^\infty e^{-r^2u^2}w(u)\,du$, avec $w\ge0$
(exemple~\ref{appA:ex:laplace}). On choisit un paramètre $\eta>0$ et l'on écrit
\[V(r)=\underbrace{\int_\eta^\infty e^{-r^2u^2}w(u)\,du}_{\text{courte portée}}+
\underbrace{\int_0^\eta e^{-r^2u^2}w(u)\,du}_{\text{lisse, longue portée}} .\]
\begin{itemize}
\item La partie de courte portée décroît comme $e^{-\eta^2r^2}$ : sa somme sur le réseau converge très vite dans
l'espace direct.
\item La partie lisse est un mélange de gaussiennes larges. Le théorème~\ref{appA:thm:poissonres} transforme sa somme
en une somme sur le réseau dual de gaussiennes étroites $e^{-\pi^2|\xi|^2/u^2}$, qui converge aussi très vite.
\end{itemize}
Le résultat ne dépend pas de $\eta$ : c'est un contrôle naturel du code. Le cas U1 du chapitre~5 rapporte
justement que l'énergie d'Ewald de rusty-SUNDIALS ne varie pas avec ce paramètre de découpage à mieux que
$9{,}0\times10^{-15}$ près \tierX.

%=====================================================================================================================
\section{Fonctions spéciales}\label{appA:sec:speciales}

Nous ne démontrons ici que les propriétés élémentaires. Les autres sont citées du \emph{Digital Library of
Mathematical Functions} du NIST~\cite{app:dlmf} (chapitres~5 pour $\Gamma$, 7 pour $\operatorname{erf}$, 10 pour les
fonctions de Bessel, 11 pour celles de Struve) \tierL.

\subsection{La fonction \texorpdfstring{$\Gamma$}{Gamma}}
Pour $s>0$, $\Gamma(s)=\int_0^\infty u^{s-1}e^{-u}du$. Une intégration par parties donne $\Gamma(s+1)=s\Gamma(s)$, donc
$\Gamma(n+1)=n!$. Le changement $u=\pi x^2$ donne $\Gamma(1/2)=\int_\R e^{-\pi x^2}dx\cdot\sqrt\pi=\sqrt\pi$, grâce à
l'intégrale de Gauss.

\subsection{Fonction d'erreur et fonction d'erreur complémentaire}
\begin{definition}
$\operatorname{erf}(x)=\frac2{\sqrt\pi}\int_0^xe^{-s^2}ds$, $\operatorname{erfc}(x)=1-\operatorname{erf}(x)=
\frac2{\sqrt\pi}\int_x^\infty e^{-s^2}ds$, et $\operatorname{erfcx}(x)=e^{x^2}\operatorname{erfc}(x)$.
\end{definition}

\begin{proposition}\label{appA:prop:erfc}
Pour $x>0$ : $\displaystyle\frac{2}{\sqrt\pi}\,\frac{e^{-x^2}}{x+\sqrt{x^2+2}}\le\operatorname{erfc}(x)
\le\frac{e^{-x^2}}{x\sqrt\pi}$. En particulier $\operatorname{erfcx}(x)\le\frac1{x\sqrt\pi}$ et
$\operatorname{erfc}(x)\sim\frac{e^{-x^2}}{x\sqrt\pi}$ quand $x\to\infty$.
\end{proposition}

\begin{humanproof}
\emph{Majoration.} Pour $s\ge x>0$, $1\le s/x$, donc $\int_x^\infty e^{-s^2}ds\le\frac1x\int_x^\infty se^{-s^2}ds=
\frac{e^{-x^2}}{2x}$. On multiplie par $2/\sqrt\pi$.

\emph{Minoration.} On l'admet \tierL{} (elle se démontre en étudiant la différence des deux membres, qui tend vers $0$ à
l'infini et dont on contrôle le signe de la dérivée) ; elle n'est utilisée que pour l'équivalent, que la majoration et
une intégration par parties suffisent aussi à établir (exercice~\ref{appA:ex:erfc}).
\end{humanproof}
Le chapitre~2 utilise les bornes $0<\operatorname{erfcx}(x)\le1$ pour $x\ge0$ et $\operatorname{erfcx}(x)\le1/(\sqrt\pi
x)$ pour $x>0$, démontrées là par la représentation $\operatorname{erfcx}(x)=\frac2{\sqrt\pi}\int_0^\infty
e^{-(y^2+2xy)}dy$. Dans les sommes de réseau certifiées, $\operatorname{erfc}$ apparaît aussi dans les bornes de queue
des gaussiennes : $\sum_{|n|>R}e^{-\pi n^2}$ se compare à une intégrale de gaussienne, c'est-à-dire à
$\operatorname{erfc}(\sqrt\pi R)$.

\subsection{Les fonctions de Bessel \texorpdfstring{$J_0$}{J0}, \texorpdfstring{$Y_0$}{Y0} et \texorpdfstring{$K_0$}{K0}}

\begin{proposition}[$J_0$]\label{appA:prop:J0}
La fonction $J_0$ définie par~\eqref{appA:eq:J0int} vérifie :
\begin{enumerate}
\item $J_0(z)=\sum_{m\ge0}\frac{(-1)^m}{(m!)^2}\bigl(\frac z2\bigr)^{2m}$ pour tout $z\in\R$ ;
\item $|J_0(z)|\le1$ et $J_0(0)=1$ ;
\item $J_0$ est solution de l'équation de Bessel d'ordre $0$ : $zy''+y'+zy=0$ ;
\item \tierL{} $J_0(z)=\sqrt{2/(\pi z)}\,\bigl(\cos(z-\pi/4)+O(1/z)\bigr)$ quand $z\to+\infty$.
\end{enumerate}
\end{proposition}

\begin{humanproof}
(1) On développe $\cos(z\cos\theta)=\sum_m(-1)^m z^{2m}\cos^{2m}\theta/(2m)!$ ; la série converge normalement en
$\theta$, et l'on intègre terme à terme. Les intégrales de Wallis donnent
$\frac1\pi\int_0^\pi\cos^{2m}\theta\,d\theta=\binom{2m}m4^{-m}$, d'où le coefficient
$\frac{1}{(2m)!}\binom{2m}{m}4^{-m}=\frac{1}{(m!)^2 4^m}$.
(2) Le module de l'intégrande de~\eqref{appA:eq:J0int} vaut $1$.
(3) Par dérivation sous l'intégrale, $J_0'(z)=-\frac1\pi\int_0^\pi\cos\theta\sin(z\cos\theta)d\theta$ et
$J_0''(z)=-\frac1\pi\int_0^\pi\cos^2\theta\cos(z\cos\theta)d\theta$. Donc
$zJ_0''+zJ_0=\frac z\pi\int_0^\pi\sin^2\theta\cos(z\cos\theta)d\theta$. Une intégration par parties, avec
$\frac{d}{d\theta}\sin(z\cos\theta)=-z\sin\theta\cos(z\cos\theta)$, transforme cette intégrale en
$\frac1\pi\int_0^\pi\cos\theta\sin(z\cos\theta)d\theta=-J_0'(z)$.
\end{humanproof}

L'équation de Bessel a une seconde solution, singulière en $0$ : la fonction de Neumann $Y_0$, avec
$Y_0(z)\sim\frac2\pi\ln z$ quand $z\to0^+$ \tierL. La \emph{fonction de Bessel modifiée} $K_0$ est définie pour $x>0$ par
\[K_0(x)=\int_0^\infty e^{-x\cosh t}\,dt .\]
Elle est positive, décroissante, se comporte comme $-\ln x$ en $0$ et comme $\sqrt{\pi/(2x)}\,e^{-x}$ à l'infini \tierL.
C'est la fonction de Green de l'opérateur $-\Delta+\kappa^2$ dans le plan : $(-\Delta+\kappa^2)\frac1{2\pi}K_0(\kappa|x|)=\delta_0$
\tierL. Les chapitres~2 et~5 notent que ce potentiel de Debye--Hückel bidimensionnel n'est \emph{pas} traité dans
ce livre, qui étudie le Yukawa tridimensionnel restreint au plan, $e^{-\kappa r}/r$.

\subsection{La fonction de Struve \texorpdfstring{$\mathbf H_0$}{H0} et le potentiel de Rytova--Keldysh}

\begin{definition}
La fonction de Struve d'ordre $0$ est $\mathbf H_0(z)=\frac2\pi\sum_{m\ge0}\frac{(-1)^mz^{2m+1}}{\bigl(1\cdot3\cdot5\cdots(2m+1)\bigr)^2}$
pour $z\in\R$.
\end{definition}
La série converge pour tout $z$ (critère de d'Alembert). On a aussi la représentation intégrale
$\mathbf H_0(z)=\frac2\pi\int_0^{\pi/2}\sin(z\cos\theta)\,d\theta$ \tierL, et $\mathbf H_0(z)-Y_0(z)\sim\frac2{\pi z}$
quand $z\to+\infty$ \tierL. Mathlib, à la version épinglée, ne connaît ni $\mathbf H_0$ ni $Y_0$ : c'est pourquoi le
chapitre~2 formalise le potentiel de Keldysh sous une forme de Laplace, et relie cette forme à la combinaison
$\mathbf H_0-Y_0$ par un argument de niveau \tierC{} qui s'appuie sur la représentation
$\mathbf K_0(z)=\mathbf H_0(z)-Y_0(z)=\frac2\pi\int_0^\infty e^{-zu}(1+u^2)^{-1/2}du$ pour $\operatorname{Re}z>0$
(DLMF~11.5.2, citée au chapitre~2~\cite{ch2dlmf}). Le comportement $\sim 2/(\pi z)$ est la loi de Coulomb $1/r$
retrouvée aux grandes distances.

%=====================================================================================================================
\section{Réseaux du plan et domaine fondamental de \texorpdfstring{$SL_2(\Z)$}{SL2(Z)}}\label{appA:sec:sl2}

\subsection{Bases d'un réseau}
Un réseau du plan est $\Lambda=\Z u\oplus\Z v$ avec $(u,v)$ base de $\R^2$ (définition~\ref{ch1:def:reseau}). Deux bases
$(u,v)$ et $(u',v')$ engendrent le même réseau si et seulement si $\binom{u'}{v'}=P\binom uv$ avec $P\in GL_2(\Z)$,
c'est-à-dire $P$ à coefficients entiers et $\det P=\pm1$.

\begin{humanproof}
Si $\Z u'\oplus\Z v'=\Z u\oplus\Z v$, chaque base s'écrit dans l'autre avec des coefficients entiers : $P$ et $P'$
entières avec $PP'=I$, donc $\det P\det P'=1$ dans $\Z$, d'où $\det P=\pm1$. Réciproquement, si $P\in GL_2(\Z)$, alors
$P^{-1}=\pm\,\mathrm{com}(P)^{\mathsf T}$ est entière, et les deux réseaux sont inclus l'un dans l'autre.
\end{humanproof}
Le covolume $|\det(u,v)|$ ne dépend donc pas de la base.

\subsection{Le paramètre de forme \texorpdfstring{$\tau$}{tau}}
Identifions $\R^2$ à $\C$. Quitte à faire une similitude directe (une rotation et une homothétie, qui ne changent pas la
« forme »), on peut supposer $u=1$. Quitte à remplacer $v$ par $-v$, on peut supposer $\operatorname{Im}v>0$. Le réseau
est alors $\Z+\Z\tau$ avec $\tau$ dans le demi-plan de Poincaré $\mathbb H=\{\operatorname{Im}\tau>0\}$. Un changement
de base direct $\binom{u'}{v'}=\begin{pmatrix}d&c\\b&a\end{pmatrix}\binom uv$ avec $ad-bc=1$ remplace $\tau$ par
\[\gamma\cdot\tau=\frac{a\tau+b}{c\tau+d},\qquad\gamma=\begin{pmatrix}a&b\\c&d\end{pmatrix}\in SL_2(\Z),\]
après renormalisation de $u'$ à $1$. Les formes de réseaux à similitude près correspondent donc aux orbites de
$SL_2(\Z)$ dans $\mathbb H$. Le calcul clé est
\begin{equation}\label{appA:eq:im}
\operatorname{Im}(\gamma\cdot\tau)=\frac{\operatorname{Im}\tau}{|c\tau+d|^2},
\end{equation}
obtenu en multipliant numérateur et dénominateur par $\overline{c\tau+d}$ et en utilisant $ad-bc=1$.

\subsection{Le domaine fondamental}
\begin{theorem}\label{appA:thm:domaine}
Soit $\mathcal F=\{\tau\in\mathbb H:|\operatorname{Re}\tau|\le\tfrac12,\ |\tau|\ge1\}$.
\begin{enumerate}
\item Tout $\tau\in\mathbb H$ a un point de son orbite sous $SL_2(\Z)$ dans $\mathcal F$.
\item Si deux points distincts de $\mathcal F$ sont dans la même orbite, ils sont sur le bord de $\mathcal F$
(symétriques par rapport à l'axe imaginaire, sur les droites $\operatorname{Re}\tau=\pm\frac12$ ou sur le cercle
$|\tau|=1$).
\end{enumerate}
\end{theorem}

\begin{humanproof}[Démonstration humaine du point 1]
Notons $T:\tau\mapsto\tau+1$ et $S:\tau\mapsto-1/\tau$ ; ils correspondent à $\bigl(\begin{smallmatrix}1&1\\0&1\end{smallmatrix}\bigr)$
et $\bigl(\begin{smallmatrix}0&-1\\1&0\end{smallmatrix}\bigr)$. Pour $\tau$ fixé, les nombres $c\tau+d$, $(c,d)\in\Z^2\setminus\{0\}$,
forment un réseau de $\C$ ; il n'y en a qu'un nombre fini de module $\le1$. D'après~\eqref{appA:eq:im},
$\operatorname{Im}(\gamma\cdot\tau)$ ne prend donc qu'un nombre fini de valeurs $\ge\operatorname{Im}\tau$, et il existe
$\gamma_0$ telle que $\tau'=\gamma_0\cdot\tau$ ait une partie imaginaire maximale dans l'orbite. Une puissance de $T$
amène $\operatorname{Re}\tau'$ dans $[-\frac12,\frac12]$ sans changer sa partie imaginaire. Si l'on avait $|\tau'|<1$,
alors $\operatorname{Im}(S\tau')=\operatorname{Im}\tau'/|\tau'|^2>\operatorname{Im}\tau'$, ce qui contredirait la
maximalité. Donc $\tau'\in\mathcal F$.
\end{humanproof}

Le point~2, dont la preuve est un examen de cas à partir de~\eqref{appA:eq:im}, est de niveau \tierL~\cite{app:serre}.
Le domaine $\mathcal F$ contient deux points remarquables :
\begin{itemize}
\item $\tau=i$, le réseau carré. Son stabilisateur dans $SL_2(\Z)$ est engendré par $S$, d'ordre $4$ dans $SL_2(\Z)$
(et $2$ dans le quotient par $\pm I$) : $S$ agit au voisinage de $i$ comme une rotation d'angle $\pi$.
\item $\tau=\rho=e^{i\pi/3}=\frac12+i\frac{\sqrt3}2$, le réseau triangulaire. Il est fixé par
$\gamma(\tau)=(\tau-1)/\tau$, d'ordre $3$ modulo $\pm I$, qui agit au voisinage de $\rho$ comme une rotation d'angle
$\pm120^\circ$.
\end{itemize}
C'est l'argument du §\ref{ch3:sec:isotropie} : une fonction de la forme du réseau, invariante par $SL_2(\Z)$ et de
classe $\mathcal C^2$, a un gradient nul en $\rho$ et une hessienne scalaire, parce que la seule matrice symétrique qui
commute avec une rotation de $120^\circ$ est un multiple de l'identité. En $i$, la rotation d'angle $\pi$ ne force que
le gradient à s'annuler : la hessienne peut y être anisotrope.

\subsection{Réduction de Gauss d'une base}
Le point~1 du théorème est aussi un algorithme. Partant d'une base $(u,v)$ :
\begin{enumerate}
\item si $|v|<|u|$, échanger $u$ et $v$ ;
\item remplacer $v$ par $v-ku$, où $k$ est l'entier le plus proche de $\operatorname{Re}(v\bar u)/|u|^2$ ;
\item recommencer tant que $|v|<|u|$.
\end{enumerate}
L'algorithme termine, car $|u|$ décroît strictement à chaque échange et ne prend qu'un nombre fini de valeurs (ce sont
des normes de vecteurs du réseau inférieures à la norme initiale). À la fin, $|u|\le|v|\le|v\pm u|$ : $u$ est un plus
court vecteur non nul et $\tau=v/u$ (au signe près) est dans $\mathcal F$. Les sommes de réseau tronquées des chapitres~3
et~5 utilisent une base ainsi réduite, pour laquelle le disque de rayon $R$ ne contient que $O(R^2)$ points de la forme
$mu+nv$ avec $|m|,|n|\le CR$.

%=====================================================================================================================
\section{Exercices}\label{appA:sec:exos}

\begin{exercise}\label{appA:ex:moments}
Démontrer par récurrence sur $m$ la formule des moments gaussiens de l'exemple~\ref{appA:ex:moments}. En déduire
$\int_\R x^{2m}e^{-x^2}dx=\Gamma(m+\tfrac12)$.
\end{exercise}

\begin{exercise}\label{appA:ex:erfc}
Montrer, par une intégration par parties, que $\operatorname{erfc}(x)=\frac{e^{-x^2}}{x\sqrt\pi}\bigl(1-\frac1{2x^2}
+O(x^{-4})\bigr)$ quand $x\to\infty$. Combien de termes de $\sum_{n\ge1}e^{-\pi n^2}$ faut-il pour une erreur inférieure
à $10^{-30}$ ?
\end{exercise}

\begin{exercise}
Appliquer le théorème~\ref{appA:thm:poisson1} à $f(x)=e^{-2\pi a|x|}$, $a>0$, dont la transformée de Fourier est
$\hat f(\xi)=\frac{a}{\pi(a^2+\xi^2)}$ (à vérifier). En déduire $\sum_{k\in\Z}\frac1{a^2+k^2}=\frac\pi a\coth(\pi a)$,
puis $\sum_{k\ge1}k^{-2}=\pi^2/6$ en faisant tendre $a$ vers $0$.
\end{exercise}

\begin{exercise}
Vérifier numériquement la proposition~\ref{appA:prop:hankel} pour $F(r)=e^{-\pi r^2}$ en évaluant $J_0$ par sa série
entière (avec \texttt{mpmath}, par exemple). Pourquoi la série entière de $J_0$ est-elle un mauvais moyen de calculer
$J_0(50)$ en virgule flottante ? (Penser à l'annulation de termes de l'ordre de $e^{50}$.)
\end{exercise}

\begin{exercise}
Réduire par l'algorithme de Gauss la base $u=(1,0)$, $v=(7{,}5;\,0{,}2)$, et donner le paramètre $\tau\in\mathcal F$
obtenu. Calculer le stabilisateur de $\rho$ dans $SL_2(\Z)$ et vérifier qu'il a six éléments.
\end{exercise}
.
% Etiquettes prefixees appA: ; references bibliographiques prefixees app: (refs_app.tex).
\chapter{Rappels d'analyse}\label{appA:chap}

Cette annexe rassemble les outils d'analyse que les chapitres utilisent sans toujours les redémontrer :
\begin{itemize}
\item l'interversion d'une limite et d'une intégrale ou d'une série (§\ref{appA:sec:param}) ;
\item les transformées de Fourier et de Laplace (§\ref{appA:sec:fourier}) ;
\item la formule sommatoire de Poisson, dans le cas général (§\ref{appA:sec:poisson}) ;
\item les fonctions spéciales rencontrées : $\Gamma$, $\operatorname{erfc}$, $J_0$, $K_0$, $Y_0$ et la fonction de
Struve $\mathbf H_0$ (§\ref{appA:sec:speciales}) ;
\item les réseaux du plan et le domaine fondamental de $SL_2(\Z)$ (§\ref{appA:sec:sl2}).
\end{itemize}
Le niveau visé est celui de la deuxième année de classe préparatoire. L'intégrale est celle de Lebesgue, mais il suffit
de savoir appliquer les deux théorèmes de convergence du programme.

\begin{warnbox}[: niveaux des énoncés de cette annexe]
Les démonstrations écrites ici sont des démonstrations humaines \tierC{} : aucune n'est vérifiée par Lean dans ce livre.
Les énoncés admis (convergence dominée, théorème de Fejér, propriétés des fonctions de Bessel et de Struve tirées du
DLMF~\cite{app:dlmf}) sont de niveau \tierL. La Mathlib épinglée contient des versions formelles de plusieurs d'entre
eux : la convergence dominée (\lean{MeasureTheory.tendsto_integral_of_dominated_convergence}, fichier
\texttt{MeasureTheory/Integral/DominatedConvergence.lean}) et la formule de Poisson (fichiers
\texttt{Analysis/Fourier/PoissonSummation.lean} et \texttt{Analysis/SpecialFunctions/Gaussian/PoissonSummation.lean}).
Nous avons seulement constaté que ces fichiers existent ; nous n'avons ni comparé leurs énoncés à ceux de cette annexe,
ni relié ces énoncés aux démonstrations ci-dessous.
\end{warnbox}

%=====================================================================================================================
\section{Séries et intégrales dépendant d'un paramètre}\label{appA:sec:param}

\subsection{Le théorème de convergence dominée}

\begin{theorem}[Convergence dominée \tierL]\label{appA:thm:cvd}
Soit $I$ un intervalle de $\R$ et $(f_n)$ une suite de fonctions continues par morceaux sur $I$. On suppose que :
\begin{enumerate}
\item $(f_n)$ converge simplement sur $I$ vers une fonction $f$ continue par morceaux ;
\item il existe $\varphi:I\to\R_+$ intégrable sur $I$ telle que $|f_n(x)|\le\varphi(x)$ pour tout $n$ et tout $x$
(hypothèse de domination).
\end{enumerate}
Alors les $f_n$ et $f$ sont intégrables sur $I$, et $\int_If_n\to\int_If$.
\end{theorem}

Le point essentiel est que la fonction dominante $\varphi$ ne dépend pas de $n$. Sans elle, la conclusion peut échouer :
$f_n=n\,\mathbf 1_{]0,1/n]}$ tend simplement vers $0$ sur $]0,1]$, alors que $\int f_n=1$ pour tout $n$.

\begin{theorem}[Interversion série--intégrale \tierL]\label{appA:thm:serieint}
Soit $(u_n)$ une suite de fonctions intégrables sur $I$ telle que $\sum_n u_n$ converge simplement vers une fonction
continue par morceaux et $\sum_n\int_I|u_n|<\infty$. Alors $\sum_nu_n$ est intégrable et
$\int_I\sum_nu_n=\sum_n\int_Iu_n$.
\end{theorem}

\subsection{Continuité et dérivation sous l'intégrale}

\begin{theorem}[Continuité sous l'intégrale]\label{appA:thm:cont}
Soit $A$ une partie de $\R$ et $f:A\times I\to\C$ telle que $f(\cdot,x)$ est continue pour tout $x$, $f(a,\cdot)$ est
continue par morceaux pour tout $a$, et $|f(a,x)|\le\varphi(x)$ avec $\varphi$ intégrable. Alors
$F(a)=\int_If(a,x)\,dx$ est continue sur $A$.
\end{theorem}

\begin{humanproof}
Soit $a_n\to a$ dans $A$. Les fonctions $x\mapsto f(a_n,x)$ convergent simplement vers $f(a,\cdot)$ et sont dominées par
$\varphi$. Le théorème~\ref{appA:thm:cvd} donne $F(a_n)\to F(a)$, et la caractérisation séquentielle conclut.
\end{humanproof}

\begin{theorem}[Dérivation sous l'intégrale]\label{appA:thm:deriv}
Soit $J$ un intervalle et $f:J\times I\to\C$. On suppose que $f(a,\cdot)$ est intégrable pour tout $a\in J$, que
$\partial f/\partial a$ existe, est continue en $a$ et continue par morceaux en $x$, et qu'il existe $\psi$ intégrable
avec $|\partial_af(a,x)|\le\psi(x)$ pour tous $a,x$. Alors $F(a)=\int_If(a,x)\,dx$ est de classe $\mathcal C^1$ sur $J$,
et $F'(a)=\int_I\partial_af(a,x)\,dx$.
\end{theorem}

\begin{humanproof}
Fixons $a\in J$ et une suite $h_n\to0$, $h_n\ne0$, avec $a+h_n\in J$. Posons
$g_n(x)=\bigl(f(a+h_n,x)-f(a,x)\bigr)/h_n$. Par définition de la dérivée partielle, $g_n(x)\to\partial_af(a,x)$ pour
tout $x$. Par l'inégalité des accroissements finis appliquée à $t\mapsto f(t,x)$ entre $a$ et $a+h_n$,
$|g_n(x)|\le\sup_t|\partial_af(t,x)|\le\psi(x)$. Le théorème~\ref{appA:thm:cvd} donne
$\bigl(F(a+h_n)-F(a)\bigr)/h_n=\int g_n\to\int\partial_af(a,x)\,dx$. La continuité de $F'$ découle du
théorème~\ref{appA:thm:cont} appliqué à $\partial_af$.
\end{humanproof}

\begin{remark}[Domination locale]
Il suffit en pratique que la domination ait lieu sur tout segment $[\alpha,\beta]\subset J$ : la dérivabilité est une
propriété locale. C'est ainsi qu'on traite les transformées de Laplace ci-dessous, où la domination n'est uniforme que
pour $t\ge\alpha>0$.
\end{remark}

\begin{example}[Moments gaussiens]\label{appA:ex:moments}
Pour $\alpha>0$, $G(\alpha)=\int_\R e^{-\pi\alpha x^2}dx=\alpha^{-1/2}$ (intégrale de Gauss). Sur $[\alpha_0,\infty[$ avec
$\alpha_0>0$, $|\partial_\alpha e^{-\pi\alpha x^2}|=\pi x^2e^{-\pi\alpha x^2}\le\pi x^2e^{-\pi\alpha_0x^2}$, qui est
intégrable. Le théorème~\ref{appA:thm:deriv} donne $\int\pi x^2e^{-\pi\alpha x^2}dx=-G'(\alpha)=\tfrac12\alpha^{-3/2}$.
En dérivant $m$ fois, $\int(\pi x^2)^me^{-\pi\alpha x^2}dx=\frac{(2m)!}{4^mm!}\,\alpha^{-1/2-m}$ (exercice~\ref{appA:ex:moments}).
C'est la même technique qui, au chapitre~1, calcule la transformée de Fourier d'une gaussienne
(théorème~\ref{ch1:thm:poisson}, étape~1).
\end{example}

\subsection{Séries de fonctions}

\begin{proposition}[Convergence normale \tierL]\label{appA:prop:normale}
Soit $(u_n)$ des fonctions continues sur une partie $A$ de $\R^d$ avec $\sum_n\sup_A|u_n|<\infty$. Alors
$\sum_nu_n$ converge uniformément sur $A$ et sa somme est continue. Si de plus les $u_n$ sont $\mathcal C^1$ sur un
ouvert et $\sum_n\sup|\nabla u_n|<\infty$, la somme est $\mathcal C^1$ et se dérive terme à terme.
\end{proposition}

\begin{example}[Séries thêta]
Pour un réseau $\Lambda$ du plan et $\alpha\ge\alpha_0>0$, on a $|x|^2\ge c\,(m^2+n^2)$ pour $x=mu+nv$ (chapitre~1),
donc $e^{-\pi\alpha|x|^2}\le e^{-\pi\alpha_0c(m^2+n^2)}$, terme général d'une série convergente. La série
$\Theta_\Lambda(\alpha)$ converge donc normalement sur $[\alpha_0,\infty[$, et elle y est $\mathcal C^\infty$ en
$\alpha$ : on dérive terme à terme autant qu'on veut. Les sommes de réseau $S(\alpha)=\sum|P|^2e^{-\pi\alpha|P|^2}$ du
chapitre~3 sont précisément $-\Theta_\Lambda'(\alpha)/\pi$.
\end{example}

%=====================================================================================================================
\section{Transformées de Fourier et de Laplace}\label{appA:sec:fourier}

\subsection{La transformée de Fourier}

\begin{definition}
Pour $f:\R^d\to\C$ intégrable, la \emph{transformée de Fourier} de $f$ est
\[\hat f(\xi)=\int_{\R^d}f(x)\,e^{-2\pi i\,\xi\cdot x}\,dx,\qquad\xi\in\R^d .\]
\end{definition}
C'est la convention du chapitre~1 : le facteur $2\pi$ est dans l'exponentielle, et la formule d'inversion n'a alors
aucun facteur devant l'intégrale. D'autres livres, en physique surtout, écrivent $\int f(x)e^{-ik\cdot x}dx$ ; on passe
de l'une à l'autre par $k=2\pi\xi$.

\begin{proposition}[Propriétés élémentaires]\label{appA:prop:fourierprop}
Soit $f$ intégrable, $a\in\R^d$, $\lambda>0$ et $M$ une matrice inversible.
\begin{enumerate}
\item $\hat f$ est continue et bornée par $\|f\|_1$.
\item Translation : $\widehat{f(\cdot-a)}(\xi)=e^{-2\pi i\,a\cdot\xi}\hat f(\xi)$.
\item Changement linéaire de variable : $\widehat{f\circ M}(\xi)=|\det M|^{-1}\hat f(M^{-\mathsf T}\xi)$.
\item Si $f$ est $\mathcal C^1$ et que $f$ et $\partial_jf$ sont intégrables, $\widehat{\partial_jf}(\xi)=2\pi
i\,\xi_j\hat f(\xi)$.
\end{enumerate}
\end{proposition}

\begin{humanproof}
(1) découle du théorème~\ref{appA:thm:cont}, avec $|f(x)e^{-2\pi i\xi\cdot x}|=|f(x)|$. (2) et (3) sont des changements
de variables $y=x-a$ et $y=Mx$ (jacobien $|\det M|$, et $\xi\cdot M^{-1}y=(M^{-\mathsf T}\xi)\cdot y$). (4) est une
intégration par parties en $x_j$ ; le terme de bord tend vers $0$ le long d'une suite de bornes, car $f$ est intégrable
et a une dérivée intégrable.
\end{humanproof}

\begin{example}[Gaussiennes]\label{appA:ex:gauss}
Pour $\alpha>0$, $\widehat{e^{-\pi\alpha|x|^2}}(\xi)=\alpha^{-d/2}e^{-\pi|\xi|^2/\alpha}$. La preuve en dimension~1
(une équation différentielle pour $\hat f$, obtenue par dérivation sous l'intégrale) est l'étape~1 de la démonstration
du théorème~\ref{ch1:thm:poisson} ; on passe à la dimension $d$ par Fubini, car l'exponentielle se factorise. En
particulier $e^{-\pi|x|^2}$ est sa propre transformée de Fourier. Pour une gaussienne de covariance quelconque,
le point~(3) de la proposition~\ref{appA:prop:fourierprop} donne la formule générale.
\end{example}

\begin{theorem}[Inversion \tierL]\label{appA:thm:inversion}
Si $f$ et $\hat f$ sont intégrables, alors $f(x)=\int\hat f(\xi)e^{2\pi i\,\xi\cdot x}d\xi$ en tout point où $f$ est
continue.
\end{theorem}

\subsection{Fonctions radiales du plan et fonction de Bessel \texorpdfstring{$J_0$}{J0}}\label{appA:ssec:hankel}

Les potentiels des chapitres~2 et~3 ne dépendent que de la distance $r=|x|$. Leur transformée de Fourier plane est
alors radiale aussi, et elle s'écrit avec la fonction de Bessel $J_0$, définie par
\begin{equation}\label{appA:eq:J0int}
J_0(z)=\frac1{2\pi}\int_0^{2\pi}e^{iz\cos\theta}\,d\theta=\frac1\pi\int_0^\pi\cos(z\cos\theta)\,d\theta .
\end{equation}

\begin{proposition}[Transformée de Hankel d'ordre $0$]\label{appA:prop:hankel}
Si $f(x)=F(|x|)$ est intégrable sur $\R^2$, alors $\hat f(\xi)=2\pi\int_0^\infty F(r)\,J_0(2\pi|\xi|r)\,r\,dr$.
\end{proposition}

\begin{humanproof}
En coordonnées polaires $x=r(\cos\theta,\sin\theta)$, et en choisissant les axes de sorte que
$\xi=(|\xi|,0)$ (l'intégrale est invariante par rotation), $\xi\cdot x=|\xi|r\cos\theta$. Par Fubini,
\[\hat f(\xi)=\int_0^\infty F(r)\Bigl(\int_0^{2\pi}e^{-2\pi i|\xi|r\cos\theta}d\theta\Bigr)r\,dr
=2\pi\int_0^\infty F(r)J_0(2\pi|\xi|r)\,r\,dr ,\]
car $J_0$ est paire (remplacer $\theta$ par $\theta+\pi$ dans~\eqref{appA:eq:J0int}).
\end{humanproof}

\begin{remark}[Contrôle sur la gaussienne]
Avec $F(r)=e^{-\pi ar^2}$, la proposition et l'exemple~\ref{appA:ex:gauss} donnent
$2\pi\int_0^\infty e^{-\pi ar^2}J_0(2\pi kr)\,r\,dr=a^{-1}e^{-\pi k^2/a}>0$ : c'est la formule employée au
chapitre~3 pour opposer la gaussienne, à transformée partout positive, au potentiel en épaulement des atomes de Rydberg
habillés, dont la transformée plane prend des valeurs négatives. Le calcul numérique de cette dernière est la
transformée de Hankel ci-dessus, évaluée par quadrature \tierX.
\end{remark}

\subsection{La transformée de Laplace}

\begin{definition}
Soit $\mu$ une mesure positive sur $[0,\infty[$ (par exemple $\mu(da)=w(a)\,da$ avec $w\ge0$). Sa \emph{transformée de
Laplace} est $\mathcal L\mu(t)=\int_0^\infty e^{-ta}\,\mu(da)$, pour les $t$ où l'intégrale converge.
\end{definition}

\begin{proposition}[Une transformée de Laplace est complètement monotone]\label{appA:prop:laplacecm}
Si $\mathcal L\mu(t)<\infty$ pour tout $t>0$, alors $\mathcal L\mu$ est $\mathcal C^\infty$ sur $]0,\infty[$ et
$(-1)^m(\mathcal L\mu)^{(m)}(t)=\int_0^\infty a^me^{-ta}\mu(da)\ge0$.
\end{proposition}

\begin{humanproof}
Sur $[t_0,\infty[$ avec $t_0>0$, on a $a^me^{-ta}\le a^me^{-t_0a}\le C_m\,e^{-t_0a/2}$, où
$C_m=\sup_{a\ge0}a^me^{-t_0a/2}<\infty$. Cette fonction est $\mu$-intégrable puisque $\mathcal L\mu(t_0/2)<\infty$.
Le théorème~\ref{appA:thm:deriv} (version mesure, même preuve) s'applique $m$ fois.
\end{humanproof}
C'est le sens « facile » du théorème de Bernstein (théorème~\ref{ch1:thm:bernstein}, proposition~\ref{ch1:prop:facile}) ;
la réciproque, toute fonction complètement monotone est une transformée de Laplace, est de niveau \tierL.

\begin{example}[Trois transformées de Laplace utilisées dans le livre]\label{appA:ex:laplace}
\begin{enumerate}
\item \emph{Puissances.} Pour $s>0$ et $t>0$, $t^{-s}=\frac1{\Gamma(s)}\int_0^\infty a^{s-1}e^{-ta}\,da$. En effet, le
changement de variable $u=ta$ donne $\int_0^\infty a^{s-1}e^{-ta}da=t^{-s}\int_0^\infty u^{s-1}e^{-u}du=t^{-s}\Gamma(s)$.
Avec $t=r^2$ et $s=1/2$, on obtient $1/r=\pi^{-1/2}\int_0^\infty a^{-1/2}e^{-ar^2}da$ : le potentiel de Coulomb est
un mélange de gaussiennes. Le changement $a=u^2$ donne la forme $1/r=\frac2{\sqrt\pi}\int_0^\infty e^{-r^2u^2}du$,
point de départ de la décomposition d'Ewald du cas U1 (chapitre~5).
\item \emph{Exciton dipolaire.} En retranchant la même formule écrite pour $r^2+d^2$,
$\frac1r-\frac1{\sqrt{r^2+d^2}}=\frac2{\sqrt\pi}\int_0^\infty e^{-r^2u^2}\bigl(1-e^{-d^2u^2}\bigr)du$, avec un poids
positif : c'est la formule du chapitre~5 (cas U1) au facteur $2$ près.
\item \emph{Subordination.} Pour $\rho>0$, $e^{-\rho\sqrt t}=\int_0^\infty e^{-ta}\,\frac{\rho}{2\sqrt\pi}\,a^{-3/2}
e^{-\rho^2/(4a)}\,da$ \tierL{} ; c'est la formule que le chapitre~2 emploie pour le potentiel de Yukawa et la
forme de Laplace du potentiel de Keldysh.
\end{enumerate}
\end{example}

%=====================================================================================================================
\section{La formule sommatoire de Poisson}\label{appA:sec:poisson}

Le chapitre~1 démontre la formule de Poisson pour les gaussiennes. Voici l'énoncé général dont elle est un cas
particulier, en dimension~$1$ d'abord, puis pour un réseau quelconque.

\begin{theorem}[Formule de Poisson en dimension 1]\label{appA:thm:poisson1}
Soit $f:\R\to\C$ continue. On suppose qu'il existe $C>0$ et $\delta>0$ tels que
\[|f(x)|\le C(1+|x|)^{-1-\delta}\quad\text{et}\quad|\hat f(\xi)|\le C(1+|\xi|)^{-1-\delta}\qquad(x,\xi\in\R).\]
Alors, pour tout $x\in\R$, $\displaystyle\sum_{n\in\Z}f(x+n)=\sum_{k\in\Z}\hat f(k)\,e^{2\pi ikx}$, les deux séries
convergeant absolument. En particulier $\sum_nf(n)=\sum_k\hat f(k)$.
\end{theorem}

\begin{humanproof}
\emph{Étape 1 : la périodisée.} Posons $F(x)=\sum_{n\in\Z}f(x+n)$. Sur $[0,1]$, $|f(x+n)|\le C(1+|n|-1)^{-1-\delta}$
pour $|n|\ge2$, terme général d'une série convergente ; la série converge donc normalement sur $[0,1]$, et par
périodicité sur $\R$. La fonction $F$ est continue (proposition~\ref{appA:prop:normale}) et $1$-périodique.

\emph{Étape 2 : ses coefficients de Fourier.} Pour $k\in\Z$, par le théorème~\ref{appA:thm:serieint} (la série des
intégrales des modules converge, car elle vaut $\int_\R|f|<\infty$),
\[c_k(F)=\int_0^1F(x)e^{-2\pi ikx}dx=\sum_{n\in\Z}\int_n^{n+1}f(y)e^{-2\pi ik(y-n)}dy=\int_\R f(y)e^{-2\pi iky}dy=\hat
f(k),\]
puisque $e^{2\pi ikn}=1$.

\emph{Étape 3 : la série de Fourier converge vers $F$.} L'hypothèse sur $\hat f$ donne $\sum_k|c_k(F)|<\infty$. La série
$G(x)=\sum_k\hat f(k)e^{2\pi ikx}$ converge donc normalement vers une fonction continue $1$-périodique, et
$c_k(G)=\hat f(k)$ (on intègre terme à terme). La fonction $H=F-G$ est continue, périodique, et tous ses coefficients
de Fourier sont nuls. Par le théorème de Fejér \tierL, les moyennes de Cesàro de la série de Fourier d'une fonction
continue périodique convergent uniformément vers elle ; ici elles sont toutes nulles, donc $H=0$.
\end{humanproof}

\begin{remark}[Pourquoi deux hypothèses de décroissance]
La décroissance de $f$ rend $F$ continue ; celle de $\hat f$ rend la série de Fourier absolument convergente. Sans la
seconde, la formule peut tomber en défaut : il existe des fonctions continues intégrables dont la série $\sum f(n)$
converge mais ne vaut pas $\sum\hat f(k)$. Les gaussiennes, et plus généralement les fonctions de Schwartz, vérifient
les deux hypothèses avec n'importe quel $\delta$.
\end{remark}

\begin{theorem}[Formule de Poisson pour un réseau]\label{appA:thm:poissonres}
Soit $\Lambda=M\Z^d$ un réseau de $\R^d$, de covolume $V=|\det M|$, et $\Lambda^*=M^{-\mathsf T}\Z^d$ son réseau dual.
Si $f$ est continue avec $|f(x)|\le C(1+|x|)^{-d-\delta}$ et $|\hat f(\xi)|\le C(1+|\xi|)^{-d-\delta}$, alors
\[\sum_{x\in\Lambda}f(x+y)=\frac1V\sum_{\xi\in\Lambda^*}\hat f(\xi)\,e^{2\pi i\,\xi\cdot y}\qquad(y\in\R^d).\]
\end{theorem}

\begin{humanproof}
On applique la version $d$-dimensionnelle de la démonstration précédente (le tore $\R^d/\Z^d$ à la place de
$\R/\Z$, et le théorème de Fejér sur le tore \tierL) à $g=f\circ M$. Les hypothèses de décroissance passent de $f$ à $g$,
car $|Mx|\ge\|M^{-1}\|^{-1}|x|$. On obtient $\sum_{n\in\Z^d}g(n+z)=\sum_{k\in\Z^d}\hat g(k)e^{2\pi ik\cdot z}$. Puis
$\hat g(k)=V^{-1}\hat f(M^{-\mathsf T}k)$ (proposition~\ref{appA:prop:fourierprop}, point~3). On pose enfin $y=Mz$ :
$k\cdot z=(M^{-\mathsf T}k)\cdot y$, et $M^{-\mathsf T}k$ parcourt $\Lambda^*$.
\end{humanproof}

Avec $f(x)=e^{-\pi\alpha|x|^2}$ et $y=0$, on retrouve exactement le théorème~\ref{ch1:thm:poisson}.

\subsection{Application : l'idée de la sommation d'Ewald}\label{appA:ssec:ewald}

Les sommes de réseau d'un potentiel à longue portée, comme $\sum_{x\in\Lambda\setminus\{0\}}|x|^{-3}$, convergent
lentement : l'erreur de troncature à la distance $R$ décroît comme $1/R$. La méthode d'Ewald coupe le potentiel en deux
grâce à une représentation en mélange de gaussiennes $V(r)=\int_0^\infty e^{-r^2u^2}w(u)\,du$, avec $w\ge0$
(exemple~\ref{appA:ex:laplace}). On choisit un paramètre $\eta>0$ et l'on écrit
\[V(r)=\underbrace{\int_\eta^\infty e^{-r^2u^2}w(u)\,du}_{\text{courte portée}}+
\underbrace{\int_0^\eta e^{-r^2u^2}w(u)\,du}_{\text{lisse, longue portée}} .\]
\begin{itemize}
\item La partie de courte portée décroît comme $e^{-\eta^2r^2}$ : sa somme sur le réseau converge très vite dans
l'espace direct.
\item La partie lisse est un mélange de gaussiennes larges. Le théorème~\ref{appA:thm:poissonres} transforme sa somme
en une somme sur le réseau dual de gaussiennes étroites $e^{-\pi^2|\xi|^2/u^2}$, qui converge aussi très vite.
\end{itemize}
Le résultat ne dépend pas de $\eta$ : c'est un contrôle naturel du code. Le cas U1 du chapitre~5 rapporte
justement que l'énergie d'Ewald de rusty-SUNDIALS ne varie pas avec ce paramètre de découpage à mieux que
$9{,}0\times10^{-15}$ près \tierX.

%=====================================================================================================================
\section{Fonctions spéciales}\label{appA:sec:speciales}

Nous ne démontrons ici que les propriétés élémentaires. Les autres sont citées du \emph{Digital Library of
Mathematical Functions} du NIST~\cite{app:dlmf} (chapitres~5 pour $\Gamma$, 7 pour $\operatorname{erf}$, 10 pour les
fonctions de Bessel, 11 pour celles de Struve) \tierL.

\subsection{La fonction \texorpdfstring{$\Gamma$}{Gamma}}
Pour $s>0$, $\Gamma(s)=\int_0^\infty u^{s-1}e^{-u}du$. Une intégration par parties donne $\Gamma(s+1)=s\Gamma(s)$, donc
$\Gamma(n+1)=n!$. Le changement $u=\pi x^2$ donne $\Gamma(1/2)=\int_\R e^{-\pi x^2}dx\cdot\sqrt\pi=\sqrt\pi$, grâce à
l'intégrale de Gauss.

\subsection{Fonction d'erreur et fonction d'erreur complémentaire}
\begin{definition}
$\operatorname{erf}(x)=\frac2{\sqrt\pi}\int_0^xe^{-s^2}ds$, $\operatorname{erfc}(x)=1-\operatorname{erf}(x)=
\frac2{\sqrt\pi}\int_x^\infty e^{-s^2}ds$, et $\operatorname{erfcx}(x)=e^{x^2}\operatorname{erfc}(x)$.
\end{definition}

\begin{proposition}\label{appA:prop:erfc}
Pour $x>0$ : $\displaystyle\frac{2}{\sqrt\pi}\,\frac{e^{-x^2}}{x+\sqrt{x^2+2}}\le\operatorname{erfc}(x)
\le\frac{e^{-x^2}}{x\sqrt\pi}$. En particulier $\operatorname{erfcx}(x)\le\frac1{x\sqrt\pi}$ et
$\operatorname{erfc}(x)\sim\frac{e^{-x^2}}{x\sqrt\pi}$ quand $x\to\infty$.
\end{proposition}

\begin{humanproof}
\emph{Majoration.} Pour $s\ge x>0$, $1\le s/x$, donc $\int_x^\infty e^{-s^2}ds\le\frac1x\int_x^\infty se^{-s^2}ds=
\frac{e^{-x^2}}{2x}$. On multiplie par $2/\sqrt\pi$.

\emph{Minoration.} On l'admet \tierL{} (elle se démontre en étudiant la différence des deux membres, qui tend vers $0$ à
l'infini et dont on contrôle le signe de la dérivée) ; elle n'est utilisée que pour l'équivalent, que la majoration et
une intégration par parties suffisent aussi à établir (exercice~\ref{appA:ex:erfc}).
\end{humanproof}
Le chapitre~2 utilise les bornes $0<\operatorname{erfcx}(x)\le1$ pour $x\ge0$ et $\operatorname{erfcx}(x)\le1/(\sqrt\pi
x)$ pour $x>0$, démontrées là par la représentation $\operatorname{erfcx}(x)=\frac2{\sqrt\pi}\int_0^\infty
e^{-(y^2+2xy)}dy$. Dans les sommes de réseau certifiées, $\operatorname{erfc}$ apparaît aussi dans les bornes de queue
des gaussiennes : $\sum_{|n|>R}e^{-\pi n^2}$ se compare à une intégrale de gaussienne, c'est-à-dire à
$\operatorname{erfc}(\sqrt\pi R)$.

\subsection{Les fonctions de Bessel \texorpdfstring{$J_0$}{J0}, \texorpdfstring{$Y_0$}{Y0} et \texorpdfstring{$K_0$}{K0}}

\begin{proposition}[$J_0$]\label{appA:prop:J0}
La fonction $J_0$ définie par~\eqref{appA:eq:J0int} vérifie :
\begin{enumerate}
\item $J_0(z)=\sum_{m\ge0}\frac{(-1)^m}{(m!)^2}\bigl(\frac z2\bigr)^{2m}$ pour tout $z\in\R$ ;
\item $|J_0(z)|\le1$ et $J_0(0)=1$ ;
\item $J_0$ est solution de l'équation de Bessel d'ordre $0$ : $zy''+y'+zy=0$ ;
\item \tierL{} $J_0(z)=\sqrt{2/(\pi z)}\,\bigl(\cos(z-\pi/4)+O(1/z)\bigr)$ quand $z\to+\infty$.
\end{enumerate}
\end{proposition}

\begin{humanproof}
(1) On développe $\cos(z\cos\theta)=\sum_m(-1)^m z^{2m}\cos^{2m}\theta/(2m)!$ ; la série converge normalement en
$\theta$, et l'on intègre terme à terme. Les intégrales de Wallis donnent
$\frac1\pi\int_0^\pi\cos^{2m}\theta\,d\theta=\binom{2m}m4^{-m}$, d'où le coefficient
$\frac{1}{(2m)!}\binom{2m}{m}4^{-m}=\frac{1}{(m!)^2 4^m}$.
(2) Le module de l'intégrande de~\eqref{appA:eq:J0int} vaut $1$.
(3) Par dérivation sous l'intégrale, $J_0'(z)=-\frac1\pi\int_0^\pi\cos\theta\sin(z\cos\theta)d\theta$ et
$J_0''(z)=-\frac1\pi\int_0^\pi\cos^2\theta\cos(z\cos\theta)d\theta$. Donc
$zJ_0''+zJ_0=\frac z\pi\int_0^\pi\sin^2\theta\cos(z\cos\theta)d\theta$. Une intégration par parties, avec
$\frac{d}{d\theta}\sin(z\cos\theta)=-z\sin\theta\cos(z\cos\theta)$, transforme cette intégrale en
$\frac1\pi\int_0^\pi\cos\theta\sin(z\cos\theta)d\theta=-J_0'(z)$.
\end{humanproof}

L'équation de Bessel a une seconde solution, singulière en $0$ : la fonction de Neumann $Y_0$, avec
$Y_0(z)\sim\frac2\pi\ln z$ quand $z\to0^+$ \tierL. La \emph{fonction de Bessel modifiée} $K_0$ est définie pour $x>0$ par
\[K_0(x)=\int_0^\infty e^{-x\cosh t}\,dt .\]
Elle est positive, décroissante, se comporte comme $-\ln x$ en $0$ et comme $\sqrt{\pi/(2x)}\,e^{-x}$ à l'infini \tierL.
C'est la fonction de Green de l'opérateur $-\Delta+\kappa^2$ dans le plan : $(-\Delta+\kappa^2)\frac1{2\pi}K_0(\kappa|x|)=\delta_0$
\tierL. Les chapitres~2 et~5 notent que ce potentiel de Debye--Hückel bidimensionnel n'est \emph{pas} traité dans
ce livre, qui étudie le Yukawa tridimensionnel restreint au plan, $e^{-\kappa r}/r$.

\subsection{La fonction de Struve \texorpdfstring{$\mathbf H_0$}{H0} et le potentiel de Rytova--Keldysh}

\begin{definition}
La fonction de Struve d'ordre $0$ est $\mathbf H_0(z)=\frac2\pi\sum_{m\ge0}\frac{(-1)^mz^{2m+1}}{\bigl(1\cdot3\cdot5\cdots(2m+1)\bigr)^2}$
pour $z\in\R$.
\end{definition}
La série converge pour tout $z$ (critère de d'Alembert). On a aussi la représentation intégrale
$\mathbf H_0(z)=\frac2\pi\int_0^{\pi/2}\sin(z\cos\theta)\,d\theta$ \tierL, et $\mathbf H_0(z)-Y_0(z)\sim\frac2{\pi z}$
quand $z\to+\infty$ \tierL. Mathlib, à la version épinglée, ne connaît ni $\mathbf H_0$ ni $Y_0$ : c'est pourquoi le
chapitre~2 formalise le potentiel de Keldysh sous une forme de Laplace, et relie cette forme à la combinaison
$\mathbf H_0-Y_0$ par un argument de niveau \tierC{} qui s'appuie sur la représentation
$\mathbf K_0(z)=\mathbf H_0(z)-Y_0(z)=\frac2\pi\int_0^\infty e^{-zu}(1+u^2)^{-1/2}du$ pour $\operatorname{Re}z>0$
(DLMF~11.5.2, citée au chapitre~2~\cite{ch2dlmf}). Le comportement $\sim 2/(\pi z)$ est la loi de Coulomb $1/r$
retrouvée aux grandes distances.

%=====================================================================================================================
\section{Réseaux du plan et domaine fondamental de \texorpdfstring{$SL_2(\Z)$}{SL2(Z)}}\label{appA:sec:sl2}

\subsection{Bases d'un réseau}
Un réseau du plan est $\Lambda=\Z u\oplus\Z v$ avec $(u,v)$ base de $\R^2$ (définition~\ref{ch1:def:reseau}). Deux bases
$(u,v)$ et $(u',v')$ engendrent le même réseau si et seulement si $\binom{u'}{v'}=P\binom uv$ avec $P\in GL_2(\Z)$,
c'est-à-dire $P$ à coefficients entiers et $\det P=\pm1$.

\begin{humanproof}
Si $\Z u'\oplus\Z v'=\Z u\oplus\Z v$, chaque base s'écrit dans l'autre avec des coefficients entiers : $P$ et $P'$
entières avec $PP'=I$, donc $\det P\det P'=1$ dans $\Z$, d'où $\det P=\pm1$. Réciproquement, si $P\in GL_2(\Z)$, alors
$P^{-1}=\pm\,\mathrm{com}(P)^{\mathsf T}$ est entière, et les deux réseaux sont inclus l'un dans l'autre.
\end{humanproof}
Le covolume $|\det(u,v)|$ ne dépend donc pas de la base.

\subsection{Le paramètre de forme \texorpdfstring{$\tau$}{tau}}
Identifions $\R^2$ à $\C$. Quitte à faire une similitude directe (une rotation et une homothétie, qui ne changent pas la
« forme »), on peut supposer $u=1$. Quitte à remplacer $v$ par $-v$, on peut supposer $\operatorname{Im}v>0$. Le réseau
est alors $\Z+\Z\tau$ avec $\tau$ dans le demi-plan de Poincaré $\mathbb H=\{\operatorname{Im}\tau>0\}$. Un changement
de base direct $\binom{u'}{v'}=\begin{pmatrix}d&c\\b&a\end{pmatrix}\binom uv$ avec $ad-bc=1$ remplace $\tau$ par
\[\gamma\cdot\tau=\frac{a\tau+b}{c\tau+d},\qquad\gamma=\begin{pmatrix}a&b\\c&d\end{pmatrix}\in SL_2(\Z),\]
après renormalisation de $u'$ à $1$. Les formes de réseaux à similitude près correspondent donc aux orbites de
$SL_2(\Z)$ dans $\mathbb H$. Le calcul clé est
\begin{equation}\label{appA:eq:im}
\operatorname{Im}(\gamma\cdot\tau)=\frac{\operatorname{Im}\tau}{|c\tau+d|^2},
\end{equation}
obtenu en multipliant numérateur et dénominateur par $\overline{c\tau+d}$ et en utilisant $ad-bc=1$.

\subsection{Le domaine fondamental}
\begin{theorem}\label{appA:thm:domaine}
Soit $\mathcal F=\{\tau\in\mathbb H:|\operatorname{Re}\tau|\le\tfrac12,\ |\tau|\ge1\}$.
\begin{enumerate}
\item Tout $\tau\in\mathbb H$ a un point de son orbite sous $SL_2(\Z)$ dans $\mathcal F$.
\item Si deux points distincts de $\mathcal F$ sont dans la même orbite, ils sont sur le bord de $\mathcal F$
(symétriques par rapport à l'axe imaginaire, sur les droites $\operatorname{Re}\tau=\pm\frac12$ ou sur le cercle
$|\tau|=1$).
\end{enumerate}
\end{theorem}

\begin{humanproof}[Démonstration humaine du point 1]
Notons $T:\tau\mapsto\tau+1$ et $S:\tau\mapsto-1/\tau$ ; ils correspondent à $\bigl(\begin{smallmatrix}1&1\\0&1\end{smallmatrix}\bigr)$
et $\bigl(\begin{smallmatrix}0&-1\\1&0\end{smallmatrix}\bigr)$. Pour $\tau$ fixé, les nombres $c\tau+d$, $(c,d)\in\Z^2\setminus\{0\}$,
forment un réseau de $\C$ ; il n'y en a qu'un nombre fini de module $\le1$. D'après~\eqref{appA:eq:im},
$\operatorname{Im}(\gamma\cdot\tau)$ ne prend donc qu'un nombre fini de valeurs $\ge\operatorname{Im}\tau$, et il existe
$\gamma_0$ telle que $\tau'=\gamma_0\cdot\tau$ ait une partie imaginaire maximale dans l'orbite. Une puissance de $T$
amène $\operatorname{Re}\tau'$ dans $[-\frac12,\frac12]$ sans changer sa partie imaginaire. Si l'on avait $|\tau'|<1$,
alors $\operatorname{Im}(S\tau')=\operatorname{Im}\tau'/|\tau'|^2>\operatorname{Im}\tau'$, ce qui contredirait la
maximalité. Donc $\tau'\in\mathcal F$.
\end{humanproof}

Le point~2, dont la preuve est un examen de cas à partir de~\eqref{appA:eq:im}, est de niveau \tierL~\cite{app:serre}.
Le domaine $\mathcal F$ contient deux points remarquables :
\begin{itemize}
\item $\tau=i$, le réseau carré. Son stabilisateur dans $SL_2(\Z)$ est engendré par $S$, d'ordre $4$ dans $SL_2(\Z)$
(et $2$ dans le quotient par $\pm I$) : $S$ agit au voisinage de $i$ comme une rotation d'angle $\pi$.
\item $\tau=\rho=e^{i\pi/3}=\frac12+i\frac{\sqrt3}2$, le réseau triangulaire. Il est fixé par
$\gamma(\tau)=(\tau-1)/\tau$, d'ordre $3$ modulo $\pm I$, qui agit au voisinage de $\rho$ comme une rotation d'angle
$\pm120^\circ$.
\end{itemize}
C'est l'argument du §\ref{ch3:sec:isotropie} : une fonction de la forme du réseau, invariante par $SL_2(\Z)$ et de
classe $\mathcal C^2$, a un gradient nul en $\rho$ et une hessienne scalaire, parce que la seule matrice symétrique qui
commute avec une rotation de $120^\circ$ est un multiple de l'identité. En $i$, la rotation d'angle $\pi$ ne force que
le gradient à s'annuler : la hessienne peut y être anisotrope.

\subsection{Réduction de Gauss d'une base}
Le point~1 du théorème est aussi un algorithme. Partant d'une base $(u,v)$ :
\begin{enumerate}
\item si $|v|<|u|$, échanger $u$ et $v$ ;
\item remplacer $v$ par $v-ku$, où $k$ est l'entier le plus proche de $\operatorname{Re}(v\bar u)/|u|^2$ ;
\item recommencer tant que $|v|<|u|$.
\end{enumerate}
L'algorithme termine, car $|u|$ décroît strictement à chaque échange et ne prend qu'un nombre fini de valeurs (ce sont
des normes de vecteurs du réseau inférieures à la norme initiale). À la fin, $|u|\le|v|\le|v\pm u|$ : $u$ est un plus
court vecteur non nul et $\tau=v/u$ (au signe près) est dans $\mathcal F$. Les sommes de réseau tronquées des chapitres~3
et~5 utilisent une base ainsi réduite, pour laquelle le disque de rayon $R$ ne contient que $O(R^2)$ points de la forme
$mu+nv$ avec $|m|,|n|\le CR$.

%=====================================================================================================================
\section{Exercices}\label{appA:sec:exos}

\begin{exercise}\label{appA:ex:moments}
Démontrer par récurrence sur $m$ la formule des moments gaussiens de l'exemple~\ref{appA:ex:moments}. En déduire
$\int_\R x^{2m}e^{-x^2}dx=\Gamma(m+\tfrac12)$.
\end{exercise}

\begin{exercise}\label{appA:ex:erfc}
Montrer, par une intégration par parties, que $\operatorname{erfc}(x)=\frac{e^{-x^2}}{x\sqrt\pi}\bigl(1-\frac1{2x^2}
+O(x^{-4})\bigr)$ quand $x\to\infty$. Combien de termes de $\sum_{n\ge1}e^{-\pi n^2}$ faut-il pour une erreur inférieure
à $10^{-30}$ ?
\end{exercise}

\begin{exercise}
Appliquer le théorème~\ref{appA:thm:poisson1} à $f(x)=e^{-2\pi a|x|}$, $a>0$, dont la transformée de Fourier est
$\hat f(\xi)=\frac{a}{\pi(a^2+\xi^2)}$ (à vérifier). En déduire $\sum_{k\in\Z}\frac1{a^2+k^2}=\frac\pi a\coth(\pi a)$,
puis $\sum_{k\ge1}k^{-2}=\pi^2/6$ en faisant tendre $a$ vers $0$.
\end{exercise}

\begin{exercise}
Vérifier numériquement la proposition~\ref{appA:prop:hankel} pour $F(r)=e^{-\pi r^2}$ en évaluant $J_0$ par sa série
entière (avec \texttt{mpmath}, par exemple). Pourquoi la série entière de $J_0$ est-elle un mauvais moyen de calculer
$J_0(50)$ en virgule flottante ? (Penser à l'annulation de termes de l'ordre de $e^{50}$.)
\end{exercise}

\begin{exercise}
Réduire par l'algorithme de Gauss la base $u=(1,0)$, $v=(7{,}5;\,0{,}2)$, et donner le paramètre $\tau\in\mathcal F$
obtenu. Calculer le stabilisateur de $\rho$ dans $SL_2(\Z)$ et vérifier qu'il a six éléments.
\end{exercise}
.
% Etiquettes prefixees appA: ; references bibliographiques prefixees app: (refs_app.tex).
\chapter{Rappels d'analyse}\label{appA:chap}

Cette annexe rassemble les outils d'analyse que les chapitres utilisent sans toujours les redémontrer :
\begin{itemize}
\item l'interversion d'une limite et d'une intégrale ou d'une série (§\ref{appA:sec:param}) ;
\item les transformées de Fourier et de Laplace (§\ref{appA:sec:fourier}) ;
\item la formule sommatoire de Poisson, dans le cas général (§\ref{appA:sec:poisson}) ;
\item les fonctions spéciales rencontrées : $\Gamma$, $\operatorname{erfc}$, $J_0$, $K_0$, $Y_0$ et la fonction de
Struve $\mathbf H_0$ (§\ref{appA:sec:speciales}) ;
\item les réseaux du plan et le domaine fondamental de $SL_2(\Z)$ (§\ref{appA:sec:sl2}).
\end{itemize}
Le niveau visé est celui de la deuxième année de classe préparatoire. L'intégrale est celle de Lebesgue, mais il suffit
de savoir appliquer les deux théorèmes de convergence du programme.

\begin{warnbox}[: niveaux des énoncés de cette annexe]
Les démonstrations écrites ici sont des démonstrations humaines \tierC{} : aucune n'est vérifiée par Lean dans ce livre.
Les énoncés admis (convergence dominée, théorème de Fejér, propriétés des fonctions de Bessel et de Struve tirées du
DLMF~\cite{app:dlmf}) sont de niveau \tierL. La Mathlib épinglée contient des versions formelles de plusieurs d'entre
eux : la convergence dominée (\lean{MeasureTheory.tendsto_integral_of_dominated_convergence}, fichier
\texttt{MeasureTheory/Integral/DominatedConvergence.lean}) et la formule de Poisson (fichiers
\texttt{Analysis/Fourier/PoissonSummation.lean} et \texttt{Analysis/SpecialFunctions/Gaussian/PoissonSummation.lean}).
Nous avons seulement constaté que ces fichiers existent ; nous n'avons ni comparé leurs énoncés à ceux de cette annexe,
ni relié ces énoncés aux démonstrations ci-dessous.
\end{warnbox}

%=====================================================================================================================
\section{Séries et intégrales dépendant d'un paramètre}\label{appA:sec:param}

\subsection{Le théorème de convergence dominée}

\begin{theorem}[Convergence dominée \tierL]\label{appA:thm:cvd}
Soit $I$ un intervalle de $\R$ et $(f_n)$ une suite de fonctions continues par morceaux sur $I$. On suppose que :
\begin{enumerate}
\item $(f_n)$ converge simplement sur $I$ vers une fonction $f$ continue par morceaux ;
\item il existe $\varphi:I\to\R_+$ intégrable sur $I$ telle que $|f_n(x)|\le\varphi(x)$ pour tout $n$ et tout $x$
(hypothèse de domination).
\end{enumerate}
Alors les $f_n$ et $f$ sont intégrables sur $I$, et $\int_If_n\to\int_If$.
\end{theorem}

Le point essentiel est que la fonction dominante $\varphi$ ne dépend pas de $n$. Sans elle, la conclusion peut échouer :
$f_n=n\,\mathbf 1_{]0,1/n]}$ tend simplement vers $0$ sur $]0,1]$, alors que $\int f_n=1$ pour tout $n$.

\begin{theorem}[Interversion série--intégrale \tierL]\label{appA:thm:serieint}
Soit $(u_n)$ une suite de fonctions intégrables sur $I$ telle que $\sum_n u_n$ converge simplement vers une fonction
continue par morceaux et $\sum_n\int_I|u_n|<\infty$. Alors $\sum_nu_n$ est intégrable et
$\int_I\sum_nu_n=\sum_n\int_Iu_n$.
\end{theorem}

\subsection{Continuité et dérivation sous l'intégrale}

\begin{theorem}[Continuité sous l'intégrale]\label{appA:thm:cont}
Soit $A$ une partie de $\R$ et $f:A\times I\to\C$ telle que $f(\cdot,x)$ est continue pour tout $x$, $f(a,\cdot)$ est
continue par morceaux pour tout $a$, et $|f(a,x)|\le\varphi(x)$ avec $\varphi$ intégrable. Alors
$F(a)=\int_If(a,x)\,dx$ est continue sur $A$.
\end{theorem}

\begin{humanproof}
Soit $a_n\to a$ dans $A$. Les fonctions $x\mapsto f(a_n,x)$ convergent simplement vers $f(a,\cdot)$ et sont dominées par
$\varphi$. Le théorème~\ref{appA:thm:cvd} donne $F(a_n)\to F(a)$, et la caractérisation séquentielle conclut.
\end{humanproof}

\begin{theorem}[Dérivation sous l'intégrale]\label{appA:thm:deriv}
Soit $J$ un intervalle et $f:J\times I\to\C$. On suppose que $f(a,\cdot)$ est intégrable pour tout $a\in J$, que
$\partial f/\partial a$ existe, est continue en $a$ et continue par morceaux en $x$, et qu'il existe $\psi$ intégrable
avec $|\partial_af(a,x)|\le\psi(x)$ pour tous $a,x$. Alors $F(a)=\int_If(a,x)\,dx$ est de classe $\mathcal C^1$ sur $J$,
et $F'(a)=\int_I\partial_af(a,x)\,dx$.
\end{theorem}

\begin{humanproof}
Fixons $a\in J$ et une suite $h_n\to0$, $h_n\ne0$, avec $a+h_n\in J$. Posons
$g_n(x)=\bigl(f(a+h_n,x)-f(a,x)\bigr)/h_n$. Par définition de la dérivée partielle, $g_n(x)\to\partial_af(a,x)$ pour
tout $x$. Par l'inégalité des accroissements finis appliquée à $t\mapsto f(t,x)$ entre $a$ et $a+h_n$,
$|g_n(x)|\le\sup_t|\partial_af(t,x)|\le\psi(x)$. Le théorème~\ref{appA:thm:cvd} donne
$\bigl(F(a+h_n)-F(a)\bigr)/h_n=\int g_n\to\int\partial_af(a,x)\,dx$. La continuité de $F'$ découle du
théorème~\ref{appA:thm:cont} appliqué à $\partial_af$.
\end{humanproof}

\begin{remark}[Domination locale]
Il suffit en pratique que la domination ait lieu sur tout segment $[\alpha,\beta]\subset J$ : la dérivabilité est une
propriété locale. C'est ainsi qu'on traite les transformées de Laplace ci-dessous, où la domination n'est uniforme que
pour $t\ge\alpha>0$.
\end{remark}

\begin{example}[Moments gaussiens]\label{appA:ex:moments}
Pour $\alpha>0$, $G(\alpha)=\int_\R e^{-\pi\alpha x^2}dx=\alpha^{-1/2}$ (intégrale de Gauss). Sur $[\alpha_0,\infty[$ avec
$\alpha_0>0$, $|\partial_\alpha e^{-\pi\alpha x^2}|=\pi x^2e^{-\pi\alpha x^2}\le\pi x^2e^{-\pi\alpha_0x^2}$, qui est
intégrable. Le théorème~\ref{appA:thm:deriv} donne $\int\pi x^2e^{-\pi\alpha x^2}dx=-G'(\alpha)=\tfrac12\alpha^{-3/2}$.
En dérivant $m$ fois, $\int(\pi x^2)^me^{-\pi\alpha x^2}dx=\frac{(2m)!}{4^mm!}\,\alpha^{-1/2-m}$ (exercice~\ref{appA:ex:moments}).
C'est la même technique qui, au chapitre~1, calcule la transformée de Fourier d'une gaussienne
(théorème~\ref{ch1:thm:poisson}, étape~1).
\end{example}

\subsection{Séries de fonctions}

\begin{proposition}[Convergence normale \tierL]\label{appA:prop:normale}
Soit $(u_n)$ des fonctions continues sur une partie $A$ de $\R^d$ avec $\sum_n\sup_A|u_n|<\infty$. Alors
$\sum_nu_n$ converge uniformément sur $A$ et sa somme est continue. Si de plus les $u_n$ sont $\mathcal C^1$ sur un
ouvert et $\sum_n\sup|\nabla u_n|<\infty$, la somme est $\mathcal C^1$ et se dérive terme à terme.
\end{proposition}

\begin{example}[Séries thêta]
Pour un réseau $\Lambda$ du plan et $\alpha\ge\alpha_0>0$, on a $|x|^2\ge c\,(m^2+n^2)$ pour $x=mu+nv$ (chapitre~1),
donc $e^{-\pi\alpha|x|^2}\le e^{-\pi\alpha_0c(m^2+n^2)}$, terme général d'une série convergente. La série
$\Theta_\Lambda(\alpha)$ converge donc normalement sur $[\alpha_0,\infty[$, et elle y est $\mathcal C^\infty$ en
$\alpha$ : on dérive terme à terme autant qu'on veut. Les sommes de réseau $S(\alpha)=\sum|P|^2e^{-\pi\alpha|P|^2}$ du
chapitre~3 sont précisément $-\Theta_\Lambda'(\alpha)/\pi$.
\end{example}

%=====================================================================================================================
\section{Transformées de Fourier et de Laplace}\label{appA:sec:fourier}

\subsection{La transformée de Fourier}

\begin{definition}
Pour $f:\R^d\to\C$ intégrable, la \emph{transformée de Fourier} de $f$ est
\[\hat f(\xi)=\int_{\R^d}f(x)\,e^{-2\pi i\,\xi\cdot x}\,dx,\qquad\xi\in\R^d .\]
\end{definition}
C'est la convention du chapitre~1 : le facteur $2\pi$ est dans l'exponentielle, et la formule d'inversion n'a alors
aucun facteur devant l'intégrale. D'autres livres, en physique surtout, écrivent $\int f(x)e^{-ik\cdot x}dx$ ; on passe
de l'une à l'autre par $k=2\pi\xi$.

\begin{proposition}[Propriétés élémentaires]\label{appA:prop:fourierprop}
Soit $f$ intégrable, $a\in\R^d$, $\lambda>0$ et $M$ une matrice inversible.
\begin{enumerate}
\item $\hat f$ est continue et bornée par $\|f\|_1$.
\item Translation : $\widehat{f(\cdot-a)}(\xi)=e^{-2\pi i\,a\cdot\xi}\hat f(\xi)$.
\item Changement linéaire de variable : $\widehat{f\circ M}(\xi)=|\det M|^{-1}\hat f(M^{-\mathsf T}\xi)$.
\item Si $f$ est $\mathcal C^1$ et que $f$ et $\partial_jf$ sont intégrables, $\widehat{\partial_jf}(\xi)=2\pi
i\,\xi_j\hat f(\xi)$.
\end{enumerate}
\end{proposition}

\begin{humanproof}
(1) découle du théorème~\ref{appA:thm:cont}, avec $|f(x)e^{-2\pi i\xi\cdot x}|=|f(x)|$. (2) et (3) sont des changements
de variables $y=x-a$ et $y=Mx$ (jacobien $|\det M|$, et $\xi\cdot M^{-1}y=(M^{-\mathsf T}\xi)\cdot y$). (4) est une
intégration par parties en $x_j$ ; le terme de bord tend vers $0$ le long d'une suite de bornes, car $f$ est intégrable
et a une dérivée intégrable.
\end{humanproof}

\begin{example}[Gaussiennes]\label{appA:ex:gauss}
Pour $\alpha>0$, $\widehat{e^{-\pi\alpha|x|^2}}(\xi)=\alpha^{-d/2}e^{-\pi|\xi|^2/\alpha}$. La preuve en dimension~1
(une équation différentielle pour $\hat f$, obtenue par dérivation sous l'intégrale) est l'étape~1 de la démonstration
du théorème~\ref{ch1:thm:poisson} ; on passe à la dimension $d$ par Fubini, car l'exponentielle se factorise. En
particulier $e^{-\pi|x|^2}$ est sa propre transformée de Fourier. Pour une gaussienne de covariance quelconque,
le point~(3) de la proposition~\ref{appA:prop:fourierprop} donne la formule générale.
\end{example}

\begin{theorem}[Inversion \tierL]\label{appA:thm:inversion}
Si $f$ et $\hat f$ sont intégrables, alors $f(x)=\int\hat f(\xi)e^{2\pi i\,\xi\cdot x}d\xi$ en tout point où $f$ est
continue.
\end{theorem}

\subsection{Fonctions radiales du plan et fonction de Bessel \texorpdfstring{$J_0$}{J0}}\label{appA:ssec:hankel}

Les potentiels des chapitres~2 et~3 ne dépendent que de la distance $r=|x|$. Leur transformée de Fourier plane est
alors radiale aussi, et elle s'écrit avec la fonction de Bessel $J_0$, définie par
\begin{equation}\label{appA:eq:J0int}
J_0(z)=\frac1{2\pi}\int_0^{2\pi}e^{iz\cos\theta}\,d\theta=\frac1\pi\int_0^\pi\cos(z\cos\theta)\,d\theta .
\end{equation}

\begin{proposition}[Transformée de Hankel d'ordre $0$]\label{appA:prop:hankel}
Si $f(x)=F(|x|)$ est intégrable sur $\R^2$, alors $\hat f(\xi)=2\pi\int_0^\infty F(r)\,J_0(2\pi|\xi|r)\,r\,dr$.
\end{proposition}

\begin{humanproof}
En coordonnées polaires $x=r(\cos\theta,\sin\theta)$, et en choisissant les axes de sorte que
$\xi=(|\xi|,0)$ (l'intégrale est invariante par rotation), $\xi\cdot x=|\xi|r\cos\theta$. Par Fubini,
\[\hat f(\xi)=\int_0^\infty F(r)\Bigl(\int_0^{2\pi}e^{-2\pi i|\xi|r\cos\theta}d\theta\Bigr)r\,dr
=2\pi\int_0^\infty F(r)J_0(2\pi|\xi|r)\,r\,dr ,\]
car $J_0$ est paire (remplacer $\theta$ par $\theta+\pi$ dans~\eqref{appA:eq:J0int}).
\end{humanproof}

\begin{remark}[Contrôle sur la gaussienne]
Avec $F(r)=e^{-\pi ar^2}$, la proposition et l'exemple~\ref{appA:ex:gauss} donnent
$2\pi\int_0^\infty e^{-\pi ar^2}J_0(2\pi kr)\,r\,dr=a^{-1}e^{-\pi k^2/a}>0$ : c'est la formule employée au
chapitre~3 pour opposer la gaussienne, à transformée partout positive, au potentiel en épaulement des atomes de Rydberg
habillés, dont la transformée plane prend des valeurs négatives. Le calcul numérique de cette dernière est la
transformée de Hankel ci-dessus, évaluée par quadrature \tierX.
\end{remark}

\subsection{La transformée de Laplace}

\begin{definition}
Soit $\mu$ une mesure positive sur $[0,\infty[$ (par exemple $\mu(da)=w(a)\,da$ avec $w\ge0$). Sa \emph{transformée de
Laplace} est $\mathcal L\mu(t)=\int_0^\infty e^{-ta}\,\mu(da)$, pour les $t$ où l'intégrale converge.
\end{definition}

\begin{proposition}[Une transformée de Laplace est complètement monotone]\label{appA:prop:laplacecm}
Si $\mathcal L\mu(t)<\infty$ pour tout $t>0$, alors $\mathcal L\mu$ est $\mathcal C^\infty$ sur $]0,\infty[$ et
$(-1)^m(\mathcal L\mu)^{(m)}(t)=\int_0^\infty a^me^{-ta}\mu(da)\ge0$.
\end{proposition}

\begin{humanproof}
Sur $[t_0,\infty[$ avec $t_0>0$, on a $a^me^{-ta}\le a^me^{-t_0a}\le C_m\,e^{-t_0a/2}$, où
$C_m=\sup_{a\ge0}a^me^{-t_0a/2}<\infty$. Cette fonction est $\mu$-intégrable puisque $\mathcal L\mu(t_0/2)<\infty$.
Le théorème~\ref{appA:thm:deriv} (version mesure, même preuve) s'applique $m$ fois.
\end{humanproof}
C'est le sens « facile » du théorème de Bernstein (théorème~\ref{ch1:thm:bernstein}, proposition~\ref{ch1:prop:facile}) ;
la réciproque, toute fonction complètement monotone est une transformée de Laplace, est de niveau \tierL.

\begin{example}[Trois transformées de Laplace utilisées dans le livre]\label{appA:ex:laplace}
\begin{enumerate}
\item \emph{Puissances.} Pour $s>0$ et $t>0$, $t^{-s}=\frac1{\Gamma(s)}\int_0^\infty a^{s-1}e^{-ta}\,da$. En effet, le
changement de variable $u=ta$ donne $\int_0^\infty a^{s-1}e^{-ta}da=t^{-s}\int_0^\infty u^{s-1}e^{-u}du=t^{-s}\Gamma(s)$.
Avec $t=r^2$ et $s=1/2$, on obtient $1/r=\pi^{-1/2}\int_0^\infty a^{-1/2}e^{-ar^2}da$ : le potentiel de Coulomb est
un mélange de gaussiennes. Le changement $a=u^2$ donne la forme $1/r=\frac2{\sqrt\pi}\int_0^\infty e^{-r^2u^2}du$,
point de départ de la décomposition d'Ewald du cas U1 (chapitre~5).
\item \emph{Exciton dipolaire.} En retranchant la même formule écrite pour $r^2+d^2$,
$\frac1r-\frac1{\sqrt{r^2+d^2}}=\frac2{\sqrt\pi}\int_0^\infty e^{-r^2u^2}\bigl(1-e^{-d^2u^2}\bigr)du$, avec un poids
positif : c'est la formule du chapitre~5 (cas U1) au facteur $2$ près.
\item \emph{Subordination.} Pour $\rho>0$, $e^{-\rho\sqrt t}=\int_0^\infty e^{-ta}\,\frac{\rho}{2\sqrt\pi}\,a^{-3/2}
e^{-\rho^2/(4a)}\,da$ \tierL{} ; c'est la formule que le chapitre~2 emploie pour le potentiel de Yukawa et la
forme de Laplace du potentiel de Keldysh.
\end{enumerate}
\end{example}

%=====================================================================================================================
\section{La formule sommatoire de Poisson}\label{appA:sec:poisson}

Le chapitre~1 démontre la formule de Poisson pour les gaussiennes. Voici l'énoncé général dont elle est un cas
particulier, en dimension~$1$ d'abord, puis pour un réseau quelconque.

\begin{theorem}[Formule de Poisson en dimension 1]\label{appA:thm:poisson1}
Soit $f:\R\to\C$ continue. On suppose qu'il existe $C>0$ et $\delta>0$ tels que
\[|f(x)|\le C(1+|x|)^{-1-\delta}\quad\text{et}\quad|\hat f(\xi)|\le C(1+|\xi|)^{-1-\delta}\qquad(x,\xi\in\R).\]
Alors, pour tout $x\in\R$, $\displaystyle\sum_{n\in\Z}f(x+n)=\sum_{k\in\Z}\hat f(k)\,e^{2\pi ikx}$, les deux séries
convergeant absolument. En particulier $\sum_nf(n)=\sum_k\hat f(k)$.
\end{theorem}

\begin{humanproof}
\emph{Étape 1 : la périodisée.} Posons $F(x)=\sum_{n\in\Z}f(x+n)$. Sur $[0,1]$, $|f(x+n)|\le C(1+|n|-1)^{-1-\delta}$
pour $|n|\ge2$, terme général d'une série convergente ; la série converge donc normalement sur $[0,1]$, et par
périodicité sur $\R$. La fonction $F$ est continue (proposition~\ref{appA:prop:normale}) et $1$-périodique.

\emph{Étape 2 : ses coefficients de Fourier.} Pour $k\in\Z$, par le théorème~\ref{appA:thm:serieint} (la série des
intégrales des modules converge, car elle vaut $\int_\R|f|<\infty$),
\[c_k(F)=\int_0^1F(x)e^{-2\pi ikx}dx=\sum_{n\in\Z}\int_n^{n+1}f(y)e^{-2\pi ik(y-n)}dy=\int_\R f(y)e^{-2\pi iky}dy=\hat
f(k),\]
puisque $e^{2\pi ikn}=1$.

\emph{Étape 3 : la série de Fourier converge vers $F$.} L'hypothèse sur $\hat f$ donne $\sum_k|c_k(F)|<\infty$. La série
$G(x)=\sum_k\hat f(k)e^{2\pi ikx}$ converge donc normalement vers une fonction continue $1$-périodique, et
$c_k(G)=\hat f(k)$ (on intègre terme à terme). La fonction $H=F-G$ est continue, périodique, et tous ses coefficients
de Fourier sont nuls. Par le théorème de Fejér \tierL, les moyennes de Cesàro de la série de Fourier d'une fonction
continue périodique convergent uniformément vers elle ; ici elles sont toutes nulles, donc $H=0$.
\end{humanproof}

\begin{remark}[Pourquoi deux hypothèses de décroissance]
La décroissance de $f$ rend $F$ continue ; celle de $\hat f$ rend la série de Fourier absolument convergente. Sans la
seconde, la formule peut tomber en défaut : il existe des fonctions continues intégrables dont la série $\sum f(n)$
converge mais ne vaut pas $\sum\hat f(k)$. Les gaussiennes, et plus généralement les fonctions de Schwartz, vérifient
les deux hypothèses avec n'importe quel $\delta$.
\end{remark}

\begin{theorem}[Formule de Poisson pour un réseau]\label{appA:thm:poissonres}
Soit $\Lambda=M\Z^d$ un réseau de $\R^d$, de covolume $V=|\det M|$, et $\Lambda^*=M^{-\mathsf T}\Z^d$ son réseau dual.
Si $f$ est continue avec $|f(x)|\le C(1+|x|)^{-d-\delta}$ et $|\hat f(\xi)|\le C(1+|\xi|)^{-d-\delta}$, alors
\[\sum_{x\in\Lambda}f(x+y)=\frac1V\sum_{\xi\in\Lambda^*}\hat f(\xi)\,e^{2\pi i\,\xi\cdot y}\qquad(y\in\R^d).\]
\end{theorem}

\begin{humanproof}
On applique la version $d$-dimensionnelle de la démonstration précédente (le tore $\R^d/\Z^d$ à la place de
$\R/\Z$, et le théorème de Fejér sur le tore \tierL) à $g=f\circ M$. Les hypothèses de décroissance passent de $f$ à $g$,
car $|Mx|\ge\|M^{-1}\|^{-1}|x|$. On obtient $\sum_{n\in\Z^d}g(n+z)=\sum_{k\in\Z^d}\hat g(k)e^{2\pi ik\cdot z}$. Puis
$\hat g(k)=V^{-1}\hat f(M^{-\mathsf T}k)$ (proposition~\ref{appA:prop:fourierprop}, point~3). On pose enfin $y=Mz$ :
$k\cdot z=(M^{-\mathsf T}k)\cdot y$, et $M^{-\mathsf T}k$ parcourt $\Lambda^*$.
\end{humanproof}

Avec $f(x)=e^{-\pi\alpha|x|^2}$ et $y=0$, on retrouve exactement le théorème~\ref{ch1:thm:poisson}.

\subsection{Application : l'idée de la sommation d'Ewald}\label{appA:ssec:ewald}

Les sommes de réseau d'un potentiel à longue portée, comme $\sum_{x\in\Lambda\setminus\{0\}}|x|^{-3}$, convergent
lentement : l'erreur de troncature à la distance $R$ décroît comme $1/R$. La méthode d'Ewald coupe le potentiel en deux
grâce à une représentation en mélange de gaussiennes $V(r)=\int_0^\infty e^{-r^2u^2}w(u)\,du$, avec $w\ge0$
(exemple~\ref{appA:ex:laplace}). On choisit un paramètre $\eta>0$ et l'on écrit
\[V(r)=\underbrace{\int_\eta^\infty e^{-r^2u^2}w(u)\,du}_{\text{courte portée}}+
\underbrace{\int_0^\eta e^{-r^2u^2}w(u)\,du}_{\text{lisse, longue portée}} .\]
\begin{itemize}
\item La partie de courte portée décroît comme $e^{-\eta^2r^2}$ : sa somme sur le réseau converge très vite dans
l'espace direct.
\item La partie lisse est un mélange de gaussiennes larges. Le théorème~\ref{appA:thm:poissonres} transforme sa somme
en une somme sur le réseau dual de gaussiennes étroites $e^{-\pi^2|\xi|^2/u^2}$, qui converge aussi très vite.
\end{itemize}
Le résultat ne dépend pas de $\eta$ : c'est un contrôle naturel du code. Le cas U1 du chapitre~5 rapporte
justement que l'énergie d'Ewald de rusty-SUNDIALS ne varie pas avec ce paramètre de découpage à mieux que
$9{,}0\times10^{-15}$ près \tierX.

%=====================================================================================================================
\section{Fonctions spéciales}\label{appA:sec:speciales}

Nous ne démontrons ici que les propriétés élémentaires. Les autres sont citées du \emph{Digital Library of
Mathematical Functions} du NIST~\cite{app:dlmf} (chapitres~5 pour $\Gamma$, 7 pour $\operatorname{erf}$, 10 pour les
fonctions de Bessel, 11 pour celles de Struve) \tierL.

\subsection{La fonction \texorpdfstring{$\Gamma$}{Gamma}}
Pour $s>0$, $\Gamma(s)=\int_0^\infty u^{s-1}e^{-u}du$. Une intégration par parties donne $\Gamma(s+1)=s\Gamma(s)$, donc
$\Gamma(n+1)=n!$. Le changement $u=\pi x^2$ donne $\Gamma(1/2)=\int_\R e^{-\pi x^2}dx\cdot\sqrt\pi=\sqrt\pi$, grâce à
l'intégrale de Gauss.

\subsection{Fonction d'erreur et fonction d'erreur complémentaire}
\begin{definition}
$\operatorname{erf}(x)=\frac2{\sqrt\pi}\int_0^xe^{-s^2}ds$, $\operatorname{erfc}(x)=1-\operatorname{erf}(x)=
\frac2{\sqrt\pi}\int_x^\infty e^{-s^2}ds$, et $\operatorname{erfcx}(x)=e^{x^2}\operatorname{erfc}(x)$.
\end{definition}

\begin{proposition}\label{appA:prop:erfc}
Pour $x>0$ : $\displaystyle\frac{2}{\sqrt\pi}\,\frac{e^{-x^2}}{x+\sqrt{x^2+2}}\le\operatorname{erfc}(x)
\le\frac{e^{-x^2}}{x\sqrt\pi}$. En particulier $\operatorname{erfcx}(x)\le\frac1{x\sqrt\pi}$ et
$\operatorname{erfc}(x)\sim\frac{e^{-x^2}}{x\sqrt\pi}$ quand $x\to\infty$.
\end{proposition}

\begin{humanproof}
\emph{Majoration.} Pour $s\ge x>0$, $1\le s/x$, donc $\int_x^\infty e^{-s^2}ds\le\frac1x\int_x^\infty se^{-s^2}ds=
\frac{e^{-x^2}}{2x}$. On multiplie par $2/\sqrt\pi$.

\emph{Minoration.} On l'admet \tierL{} (elle se démontre en étudiant la différence des deux membres, qui tend vers $0$ à
l'infini et dont on contrôle le signe de la dérivée) ; elle n'est utilisée que pour l'équivalent, que la majoration et
une intégration par parties suffisent aussi à établir (exercice~\ref{appA:ex:erfc}).
\end{humanproof}
Le chapitre~2 utilise les bornes $0<\operatorname{erfcx}(x)\le1$ pour $x\ge0$ et $\operatorname{erfcx}(x)\le1/(\sqrt\pi
x)$ pour $x>0$, démontrées là par la représentation $\operatorname{erfcx}(x)=\frac2{\sqrt\pi}\int_0^\infty
e^{-(y^2+2xy)}dy$. Dans les sommes de réseau certifiées, $\operatorname{erfc}$ apparaît aussi dans les bornes de queue
des gaussiennes : $\sum_{|n|>R}e^{-\pi n^2}$ se compare à une intégrale de gaussienne, c'est-à-dire à
$\operatorname{erfc}(\sqrt\pi R)$.

\subsection{Les fonctions de Bessel \texorpdfstring{$J_0$}{J0}, \texorpdfstring{$Y_0$}{Y0} et \texorpdfstring{$K_0$}{K0}}

\begin{proposition}[$J_0$]\label{appA:prop:J0}
La fonction $J_0$ définie par~\eqref{appA:eq:J0int} vérifie :
\begin{enumerate}
\item $J_0(z)=\sum_{m\ge0}\frac{(-1)^m}{(m!)^2}\bigl(\frac z2\bigr)^{2m}$ pour tout $z\in\R$ ;
\item $|J_0(z)|\le1$ et $J_0(0)=1$ ;
\item $J_0$ est solution de l'équation de Bessel d'ordre $0$ : $zy''+y'+zy=0$ ;
\item \tierL{} $J_0(z)=\sqrt{2/(\pi z)}\,\bigl(\cos(z-\pi/4)+O(1/z)\bigr)$ quand $z\to+\infty$.
\end{enumerate}
\end{proposition}

\begin{humanproof}
(1) On développe $\cos(z\cos\theta)=\sum_m(-1)^m z^{2m}\cos^{2m}\theta/(2m)!$ ; la série converge normalement en
$\theta$, et l'on intègre terme à terme. Les intégrales de Wallis donnent
$\frac1\pi\int_0^\pi\cos^{2m}\theta\,d\theta=\binom{2m}m4^{-m}$, d'où le coefficient
$\frac{1}{(2m)!}\binom{2m}{m}4^{-m}=\frac{1}{(m!)^2 4^m}$.
(2) Le module de l'intégrande de~\eqref{appA:eq:J0int} vaut $1$.
(3) Par dérivation sous l'intégrale, $J_0'(z)=-\frac1\pi\int_0^\pi\cos\theta\sin(z\cos\theta)d\theta$ et
$J_0''(z)=-\frac1\pi\int_0^\pi\cos^2\theta\cos(z\cos\theta)d\theta$. Donc
$zJ_0''+zJ_0=\frac z\pi\int_0^\pi\sin^2\theta\cos(z\cos\theta)d\theta$. Une intégration par parties, avec
$\frac{d}{d\theta}\sin(z\cos\theta)=-z\sin\theta\cos(z\cos\theta)$, transforme cette intégrale en
$\frac1\pi\int_0^\pi\cos\theta\sin(z\cos\theta)d\theta=-J_0'(z)$.
\end{humanproof}

L'équation de Bessel a une seconde solution, singulière en $0$ : la fonction de Neumann $Y_0$, avec
$Y_0(z)\sim\frac2\pi\ln z$ quand $z\to0^+$ \tierL. La \emph{fonction de Bessel modifiée} $K_0$ est définie pour $x>0$ par
\[K_0(x)=\int_0^\infty e^{-x\cosh t}\,dt .\]
Elle est positive, décroissante, se comporte comme $-\ln x$ en $0$ et comme $\sqrt{\pi/(2x)}\,e^{-x}$ à l'infini \tierL.
C'est la fonction de Green de l'opérateur $-\Delta+\kappa^2$ dans le plan : $(-\Delta+\kappa^2)\frac1{2\pi}K_0(\kappa|x|)=\delta_0$
\tierL. Les chapitres~2 et~5 notent que ce potentiel de Debye--Hückel bidimensionnel n'est \emph{pas} traité dans
ce livre, qui étudie le Yukawa tridimensionnel restreint au plan, $e^{-\kappa r}/r$.

\subsection{La fonction de Struve \texorpdfstring{$\mathbf H_0$}{H0} et le potentiel de Rytova--Keldysh}

\begin{definition}
La fonction de Struve d'ordre $0$ est $\mathbf H_0(z)=\frac2\pi\sum_{m\ge0}\frac{(-1)^mz^{2m+1}}{\bigl(1\cdot3\cdot5\cdots(2m+1)\bigr)^2}$
pour $z\in\R$.
\end{definition}
La série converge pour tout $z$ (critère de d'Alembert). On a aussi la représentation intégrale
$\mathbf H_0(z)=\frac2\pi\int_0^{\pi/2}\sin(z\cos\theta)\,d\theta$ \tierL, et $\mathbf H_0(z)-Y_0(z)\sim\frac2{\pi z}$
quand $z\to+\infty$ \tierL. Mathlib, à la version épinglée, ne connaît ni $\mathbf H_0$ ni $Y_0$ : c'est pourquoi le
chapitre~2 formalise le potentiel de Keldysh sous une forme de Laplace, et relie cette forme à la combinaison
$\mathbf H_0-Y_0$ par un argument de niveau \tierC{} qui s'appuie sur la représentation
$\mathbf K_0(z)=\mathbf H_0(z)-Y_0(z)=\frac2\pi\int_0^\infty e^{-zu}(1+u^2)^{-1/2}du$ pour $\operatorname{Re}z>0$
(DLMF~11.5.2, citée au chapitre~2~\cite{ch2dlmf}). Le comportement $\sim 2/(\pi z)$ est la loi de Coulomb $1/r$
retrouvée aux grandes distances.

%=====================================================================================================================
\section{Réseaux du plan et domaine fondamental de \texorpdfstring{$SL_2(\Z)$}{SL2(Z)}}\label{appA:sec:sl2}

\subsection{Bases d'un réseau}
Un réseau du plan est $\Lambda=\Z u\oplus\Z v$ avec $(u,v)$ base de $\R^2$ (définition~\ref{ch1:def:reseau}). Deux bases
$(u,v)$ et $(u',v')$ engendrent le même réseau si et seulement si $\binom{u'}{v'}=P\binom uv$ avec $P\in GL_2(\Z)$,
c'est-à-dire $P$ à coefficients entiers et $\det P=\pm1$.

\begin{humanproof}
Si $\Z u'\oplus\Z v'=\Z u\oplus\Z v$, chaque base s'écrit dans l'autre avec des coefficients entiers : $P$ et $P'$
entières avec $PP'=I$, donc $\det P\det P'=1$ dans $\Z$, d'où $\det P=\pm1$. Réciproquement, si $P\in GL_2(\Z)$, alors
$P^{-1}=\pm\,\mathrm{com}(P)^{\mathsf T}$ est entière, et les deux réseaux sont inclus l'un dans l'autre.
\end{humanproof}
Le covolume $|\det(u,v)|$ ne dépend donc pas de la base.

\subsection{Le paramètre de forme \texorpdfstring{$\tau$}{tau}}
Identifions $\R^2$ à $\C$. Quitte à faire une similitude directe (une rotation et une homothétie, qui ne changent pas la
« forme »), on peut supposer $u=1$. Quitte à remplacer $v$ par $-v$, on peut supposer $\operatorname{Im}v>0$. Le réseau
est alors $\Z+\Z\tau$ avec $\tau$ dans le demi-plan de Poincaré $\mathbb H=\{\operatorname{Im}\tau>0\}$. Un changement
de base direct $\binom{u'}{v'}=\begin{pmatrix}d&c\\b&a\end{pmatrix}\binom uv$ avec $ad-bc=1$ remplace $\tau$ par
\[\gamma\cdot\tau=\frac{a\tau+b}{c\tau+d},\qquad\gamma=\begin{pmatrix}a&b\\c&d\end{pmatrix}\in SL_2(\Z),\]
après renormalisation de $u'$ à $1$. Les formes de réseaux à similitude près correspondent donc aux orbites de
$SL_2(\Z)$ dans $\mathbb H$. Le calcul clé est
\begin{equation}\label{appA:eq:im}
\operatorname{Im}(\gamma\cdot\tau)=\frac{\operatorname{Im}\tau}{|c\tau+d|^2},
\end{equation}
obtenu en multipliant numérateur et dénominateur par $\overline{c\tau+d}$ et en utilisant $ad-bc=1$.

\subsection{Le domaine fondamental}
\begin{theorem}\label{appA:thm:domaine}
Soit $\mathcal F=\{\tau\in\mathbb H:|\operatorname{Re}\tau|\le\tfrac12,\ |\tau|\ge1\}$.
\begin{enumerate}
\item Tout $\tau\in\mathbb H$ a un point de son orbite sous $SL_2(\Z)$ dans $\mathcal F$.
\item Si deux points distincts de $\mathcal F$ sont dans la même orbite, ils sont sur le bord de $\mathcal F$
(symétriques par rapport à l'axe imaginaire, sur les droites $\operatorname{Re}\tau=\pm\frac12$ ou sur le cercle
$|\tau|=1$).
\end{enumerate}
\end{theorem}

\begin{humanproof}[Démonstration humaine du point 1]
Notons $T:\tau\mapsto\tau+1$ et $S:\tau\mapsto-1/\tau$ ; ils correspondent à $\bigl(\begin{smallmatrix}1&1\\0&1\end{smallmatrix}\bigr)$
et $\bigl(\begin{smallmatrix}0&-1\\1&0\end{smallmatrix}\bigr)$. Pour $\tau$ fixé, les nombres $c\tau+d$, $(c,d)\in\Z^2\setminus\{0\}$,
forment un réseau de $\C$ ; il n'y en a qu'un nombre fini de module $\le1$. D'après~\eqref{appA:eq:im},
$\operatorname{Im}(\gamma\cdot\tau)$ ne prend donc qu'un nombre fini de valeurs $\ge\operatorname{Im}\tau$, et il existe
$\gamma_0$ telle que $\tau'=\gamma_0\cdot\tau$ ait une partie imaginaire maximale dans l'orbite. Une puissance de $T$
amène $\operatorname{Re}\tau'$ dans $[-\frac12,\frac12]$ sans changer sa partie imaginaire. Si l'on avait $|\tau'|<1$,
alors $\operatorname{Im}(S\tau')=\operatorname{Im}\tau'/|\tau'|^2>\operatorname{Im}\tau'$, ce qui contredirait la
maximalité. Donc $\tau'\in\mathcal F$.
\end{humanproof}

Le point~2, dont la preuve est un examen de cas à partir de~\eqref{appA:eq:im}, est de niveau \tierL~\cite{app:serre}.
Le domaine $\mathcal F$ contient deux points remarquables :
\begin{itemize}
\item $\tau=i$, le réseau carré. Son stabilisateur dans $SL_2(\Z)$ est engendré par $S$, d'ordre $4$ dans $SL_2(\Z)$
(et $2$ dans le quotient par $\pm I$) : $S$ agit au voisinage de $i$ comme une rotation d'angle $\pi$.
\item $\tau=\rho=e^{i\pi/3}=\frac12+i\frac{\sqrt3}2$, le réseau triangulaire. Il est fixé par
$\gamma(\tau)=(\tau-1)/\tau$, d'ordre $3$ modulo $\pm I$, qui agit au voisinage de $\rho$ comme une rotation d'angle
$\pm120^\circ$.
\end{itemize}
C'est l'argument du §\ref{ch3:sec:isotropie} : une fonction de la forme du réseau, invariante par $SL_2(\Z)$ et de
classe $\mathcal C^2$, a un gradient nul en $\rho$ et une hessienne scalaire, parce que la seule matrice symétrique qui
commute avec une rotation de $120^\circ$ est un multiple de l'identité. En $i$, la rotation d'angle $\pi$ ne force que
le gradient à s'annuler : la hessienne peut y être anisotrope.

\subsection{Réduction de Gauss d'une base}
Le point~1 du théorème est aussi un algorithme. Partant d'une base $(u,v)$ :
\begin{enumerate}
\item si $|v|<|u|$, échanger $u$ et $v$ ;
\item remplacer $v$ par $v-ku$, où $k$ est l'entier le plus proche de $\operatorname{Re}(v\bar u)/|u|^2$ ;
\item recommencer tant que $|v|<|u|$.
\end{enumerate}
L'algorithme termine, car $|u|$ décroît strictement à chaque échange et ne prend qu'un nombre fini de valeurs (ce sont
des normes de vecteurs du réseau inférieures à la norme initiale). À la fin, $|u|\le|v|\le|v\pm u|$ : $u$ est un plus
court vecteur non nul et $\tau=v/u$ (au signe près) est dans $\mathcal F$. Les sommes de réseau tronquées des chapitres~3
et~5 utilisent une base ainsi réduite, pour laquelle le disque de rayon $R$ ne contient que $O(R^2)$ points de la forme
$mu+nv$ avec $|m|,|n|\le CR$.

%=====================================================================================================================
\section{Exercices}\label{appA:sec:exos}

\begin{exercise}\label{appA:ex:moments}
Démontrer par récurrence sur $m$ la formule des moments gaussiens de l'exemple~\ref{appA:ex:moments}. En déduire
$\int_\R x^{2m}e^{-x^2}dx=\Gamma(m+\tfrac12)$.
\end{exercise}

\begin{exercise}\label{appA:ex:erfc}
Montrer, par une intégration par parties, que $\operatorname{erfc}(x)=\frac{e^{-x^2}}{x\sqrt\pi}\bigl(1-\frac1{2x^2}
+O(x^{-4})\bigr)$ quand $x\to\infty$. Combien de termes de $\sum_{n\ge1}e^{-\pi n^2}$ faut-il pour une erreur inférieure
à $10^{-30}$ ?
\end{exercise}

\begin{exercise}
Appliquer le théorème~\ref{appA:thm:poisson1} à $f(x)=e^{-2\pi a|x|}$, $a>0$, dont la transformée de Fourier est
$\hat f(\xi)=\frac{a}{\pi(a^2+\xi^2)}$ (à vérifier). En déduire $\sum_{k\in\Z}\frac1{a^2+k^2}=\frac\pi a\coth(\pi a)$,
puis $\sum_{k\ge1}k^{-2}=\pi^2/6$ en faisant tendre $a$ vers $0$.
\end{exercise}

\begin{exercise}
Vérifier numériquement la proposition~\ref{appA:prop:hankel} pour $F(r)=e^{-\pi r^2}$ en évaluant $J_0$ par sa série
entière (avec \texttt{mpmath}, par exemple). Pourquoi la série entière de $J_0$ est-elle un mauvais moyen de calculer
$J_0(50)$ en virgule flottante ? (Penser à l'annulation de termes de l'ordre de $e^{50}$.)
\end{exercise}

\begin{exercise}
Réduire par l'algorithme de Gauss la base $u=(1,0)$, $v=(7{,}5;\,0{,}2)$, et donner le paramètre $\tau\in\mathcal F$
obtenu. Calculer le stabilisateur de $\rho$ dans $SL_2(\Z)$ et vérifier qu'il a six éléments.
\end{exercise}
